\documentclass[10pt,a4paper]{article}
% Оформление: ivanshemidko/algebra-notes, s2/main.tex,
% commit d52e89ab46b905aac035e556794af91aec7c86ef.
\usepackage{fontspec}
% Project-local installation keeps local and CI font files identical.
\setmainfont{cmunrm.ttf}[
  Path=fonts/,
  BoldFont=cmunbx.ttf,
  ItalicFont=cmunti.ttf,
  BoldItalicFont=cmunbi.ttf
]
\usepackage[english,russian]{babel}
\usepackage{amsmath,amssymb,mathtools,amsthm}
% Unicode maps for traditional Computer Modern math fonts (search/copy text).
\pdfvariable gentounicode=1
\protected\def\pdfglyphtounicode{\pdfextension glyphtounicode}
\input{glyphtounicode.tex}
% The standard map omits CMEX's sized delimiter glyph names.
\ExplSyntaxOn
\clist_map_inline:nn {big,Big,bigg,Bigg} {
  \pdfglyphtounicode{parenleft#1}{0028}
  \pdfglyphtounicode{parenright#1}{0029}
  \pdfglyphtounicode{braceleft#1}{007B}
  \pdfglyphtounicode{braceright#1}{007D}
  \pdfglyphtounicode{bracketleft#1}{005B}
  \pdfglyphtounicode{bracketright#1}{005D}
  \pdfglyphtounicode{floorleft#1}{230A}
  \pdfglyphtounicode{floorright#1}{230B}
  \pdfglyphtounicode{ceilingleft#1}{2308}
  \pdfglyphtounicode{ceilingright#1}{2309}
  \pdfglyphtounicode{angbracketleft#1}{27E8}
  \pdfglyphtounicode{angbracketright#1}{27E9}
  \pdfglyphtounicode{slash#1}{002F}
  \pdfglyphtounicode{backslash#1}{005C}
}
\ExplSyntaxOff
\usepackage{geometry}
\geometry{left=22mm,right=22mm,top=25mm,bottom=27mm,heightrounded}
\usepackage{enumitem}
\usepackage[most]{tcolorbox}
\tcbuselibrary{theorems,breakable,skins}
\usepackage{graphicx,caption}
\captionsetup{font=small,labelfont=bf,hypcap=false}

\definecolor{cbBlueBack}{RGB}{246,249,252}
\definecolor{cbGreenBack}{RGB}{246,251,248}
\definecolor{cbOrangeBack}{RGB}{252,249,242}
\definecolor{cbPurpleBack}{RGB}{249,247,252}
\definecolor{cbGrayBack}{RGB}{249,249,249}
\definecolor{cbBlue}{RGB}{42,87,132}
\definecolor{cbGreen}{RGB}{52,110,72}
\definecolor{cbOrange}{RGB}{140,98,0}
\definecolor{cbPurple}{RGB}{93,78,140}
\definecolor{cbGray}{RGB}{80,80,80}

\newtheoremstyle{mystatement}{8pt}{8pt}{\itshape}{}{\bfseries}{.}{\newline}{}
\newtheoremstyle{mydefinition}{8pt}{8pt}{\normalfont}{}{\bfseries}{.}{0.5em}{}
\theoremstyle{mystatement}
\newtheorem{theorem}{Теорема}[part]
\newtheorem{lemma}{Лемма}[part]
\newtheorem{cor}{Следствие}[part]
\theoremstyle{mydefinition}
\newtheorem{definition}{Определение}[part]
\newtheorem{example}{Пример}[part]
\numberwithin{equation}{part}
\counterwithin{figure}{part}
\tcbset{
  enhanced,sharp corners,boxrule=0.5pt,
  left=9pt,right=9pt,top=6pt,bottom=6pt,
  before skip=6pt,after skip=6pt,breakable
}
\tcolorboxenvironment{theorem}{colback=cbBlueBack,colframe=cbGray!55,borderline west={1.6pt}{0pt}{cbBlue}}
\tcolorboxenvironment{lemma}{colback=cbBlueBack,colframe=cbGray!55,borderline west={1.6pt}{0pt}{cbBlue}}
\tcolorboxenvironment{cor}{colback=cbPurpleBack,colframe=cbGray!55,borderline west={1.6pt}{0pt}{cbPurple}}
\tcolorboxenvironment{definition}{colback=cbGreenBack,colframe=cbGreen!55!cbGray,borderline west={1.6pt}{0pt}{cbGreen}}
\tcolorboxenvironment{example}{colback=cbOrangeBack,colframe=cbOrange!45!cbGray,borderline west={1.6pt}{0pt}{cbOrange}}

\usepackage{hyperref}
\hypersetup{
  colorlinks=true,linkcolor=cbBlue,citecolor=cbBlue,urlcolor=cbBlue,
  pdftitle={Дискретная математика · Коллоквиум 1},
  pdfauthor={},pdfsubject={},pdfcreator={LuaLaTeX}
}
\linespread{1.0}
\setlength{\parindent}{0pt}
\setlength{\parskip}{4pt}
\allowdisplaybreaks[1]
\providecommand{\tightlist}{\setlength{\itemsep}{0pt}\setlength{\parskip}{0pt}}

\usepackage{titlesec}
\titleformat{\part}[block]{\Large\bfseries\centering}{}{0pt}{}
\titlespacing*{\part}{0pt}{0pt}{12pt}
\counterwithin{section}{part}
\renewcommand{\thepart}{\arabic{part}}
\usepackage{tocloft}
\setlength{\cftsecnumwidth}{2.8em}
\setlength{\cftsubsecnumwidth}{3.6em}
\setcounter{tocdepth}{2}
\usepackage{fancyhdr}
\pagestyle{fancy}
\fancyhf{}
\fancyhead[L]{\nouppercase{\leftmark}}
\fancyfoot[C]{\thepage}
\let\oldpart\part
\renewcommand{\part}[1]{\oldpart{#1}\markboth{#1}{}}
\renewcommand{\sectionmark}[1]{}
\usepackage{etoolbox}
\pretocmd{\part}{\newpage}{}{}

\usepackage{tikz}
\usetikzlibrary{arrows.meta,positioning,calc,patterns,fit}
% Общий язык учебных схем; штрихи, двойной контур и метки читаются без цвета.
\tikzset{
  gi vertex/.style={circle,draw=black,fill=white,minimum size=5mm,inner sep=1pt},
  gi small/.style={circle,draw=black,fill=white,minimum size=2.5mm,inner sep=0pt},
  gi selected/.style={gi vertex,double,fill=black!12},
  gi edge/.style={draw=black,line width=.65pt},
  gi emphasis/.style={draw=cbBlue,line width=1.25pt},
  gi alternate/.style={draw=black,densely dashed,line width=1pt},
  gi removed/.style={draw=black!55,densely dotted,line width=.85pt},
  gi arrow/.style={gi edge,-{Stealth[length=2mm]}},
  gi flow/.style={-{Stealth[length=2mm]},thick},
  gi group/.style={draw=black!60,dashed,rounded corners,inner sep=6pt},
  gi note/.style={align=center,font=\small},
  gi box/.style={draw,rounded corners,align=center,inner sep=5pt,font=\small}
}
% Неплавающая композиция: рисунок и подпись переходят на страницу вместе.
\newenvironment{graphillustration}[2]{%
  \par\addvspace{6pt}\noindent\begin{minipage}{\linewidth}%
  \def\gicaption{#1}\def\gilabel{#2}\centering\normalfont\small
}{%
  \par\captionof{figure}{\gicaption}\label{\gilabel}%
  \end{minipage}\par\addvspace{6pt}%
}
\DeclareMathOperator{\comp}{c}
\DeclareMathOperator{\diag}{diag}
\DeclareMathOperator{\cl}{cl}
\DeclareMathOperator{\rank}{rank}
\newcommand{\N}{\mathbb N}
\newcommand{\Z}{\mathbb Z}
\newcommand{\R}{\mathbb R}


% Билеты используют счётчик part, но компактные строки оглавления без двойной нумерации.
\newcommand{\ticket}[1]{%
  \clearpage\refstepcounter{part}%
  \markboth{Билет \thepart. #1}{}%
  \addcontentsline{toc}{part}{Билет \thepart. #1}%
  {\centering\Large\bfseries Билет \thepart. #1\par}%
  \vspace{12pt}%
}
\makeatletter
\renewcommand*{\l@part}[2]{%
  \addpenalty{-100}\addvspace{5pt}%
  \@dottedtocline{0}{0em}{0em}{\bfseries #1}{#2}%
}
\makeatother

\begin{document}
\thispagestyle{empty}
\begin{center}
{\LARGE\bfseries Дискретная математика · III семестр}\par\medskip
{\Large Теория графов. Коллоквиум 1}
\end{center}

\clearpage
\tableofcontents
\clearpage
\phantomsection
\section*{Обозначения и соглашения}
\markboth{Обозначения и соглашения}{}
\addcontentsline{toc}{section}{Обозначения и соглашения}
\begingroup
\small
\setlength{\parskip}{2pt}
\tcbset{before skip=4pt,after skip=4pt,top=4pt,bottom=4pt}

\source{Лекции, с. 1–3, 10, 13–15, 19–20; обозначения дополнены для билетов}
\begin{definition}[Граф и подграф]
Неориентированный граф \(G=(V,E)\) состоит из конечного непустого множества
вершин \(V\) и множества неупорядоченных пар различных вершин \(E\).
Таким образом, по умолчанию петель и кратных рёбер нет.
Орграф задаётся множеством дуг \(E\subseteq V^2\); дуга \((u,v)\)
направлена из \(u\) в \(v\). Забывая направления дуг, удаляя петли и кратности,
получаем соответствующий неориентированный граф.
Подграф \(H=(W,F)\) удовлетворяет \(W\subseteq V\), \(F\subseteq E\),
причём концы всех рёбер \(F\) лежат в \(W\).
Подграф индуцирован множеством \(W\), если содержит все рёбра исходного
графа с обоими концами в \(W\); он обозначается \(G[W]\).
Остовный подграф содержит все вершины: \(W=V\).
\end{definition}
Для орграфа направление каждой дуги учитывается. Иные соглашения оговорены в билетах.
Пишем \(G=(V,E)\), \(n=|V|\), \(m=|E|\), \(c(G)\) — число компонент связности,
\(N(v)\) — множество соседей, \(d(v)=\deg(v)=|N(v)|\),
\(\delta(G)=\min d(v)\), \(\Delta(G)=\max d(v)\).
\(G-v\) и \(G-e\) — удаление вершины со всеми инцидентными рёбрами и удаление ребра;
\(G[S]\) — индуцированный подграф на \(S\).
\(K_n\), \(K_{p,q}\), \(P_n\), \(C_n\) — полный граф, полный двудольный граф,
простой путь и простой цикл соответственно; в \(P_n,C_n\) индекс означает число вершин.

\textbf{Путь} (в терминологии лекций; также \textbf{маршрут}) — последовательность \(v_0,e_1,v_1,\ldots,e_k,v_k\),
где каждое ребро соединяет соседние вершины последовательности.
Его длина равна \(k\); разрешены повторения вершин и рёбер.
В орграфе дуги проходятся по направлению.
\textbf{Цепь} — маршрут без повторений рёбер; \textbf{простой путь} — маршрут
без повторений вершин. Путь с одинаковыми началом и концом называется замкнутым.
В лекционном определении неориентированного \textbf{цикла} дополнительно
требуются положительная длина и отсутствие немедленного возвращения
по тому же ребру, в том числе на стыке конца и начала записи.
Поэтому путь \(u,v,u\) по одному ребру замкнут, но циклом не считается.
Замкнутая цепь дополнительно не повторяет рёбер;
\textbf{простой цикл} — замкнутый маршрут положительной длины без других
повторений вершин (в простом неориентированном графе длины не менее трёх).

Записи неориентированного цикла, отличающиеся циклическим сдвигом
или обращением обхода, описывают один цикл. Для ориентированного цикла
обращение допускается только при наличии нужных дуг; при подсчёте замкнутых
путей с отмеченным началом сдвиги считаются разными записями.
Цикл может проходить один и тот же цикл несколько раз; простой — нет.
Если \(P=R^r\), \(Q=R^s\) для положительных \(r,s\), где степень означает
склейку повторных обходов, то в лекциях такие циклы также отождествляются
отдельным отношением \(P\equiv Q\) (после сдвига/обращения записи).
Это отождествление не применяется при подсчёте путей заданной длины.

\textbf{Лемма (лекции, с. 15).} Наличие цикла в неориентированном графе
влечёт наличие простого цикла. Выбираем в записи цикла кратчайший
непустой замкнутый участок между повторениями вершины. Внутренних
повторений у него нет по минимальности; длина 1 исключена отсутствием
петель, длина 2 — запретом немедленного возврата по ребру.
Следовательно, участок — простой цикл.

Из любого пути между различными вершинами можно получить простой путь:
при повторении вершины вырезать участок между двумя её вхождениями.
Длина строго уменьшается, поэтому процесс заканчивается.
Граф \textbf{связен}, если любые две его вершины соединены путём.
Связность является отношением эквивалентности: нулевой путь даёт
рефлексивность, обращение пути — симметричность, склейка — транзитивность.
Её классы — компоненты связности.
Одновершинный граф связен. При удалениях допускается пустой граф:
у него ноль компонент. Это вспомогательное расширение лекционного соглашения.

\textbf{Лемма о рукопожатиях.} \(\sum_{v\in V}d(v)=2m\).
Действительно, каждое ребро учитывается ровно по одному разу у каждого конца.
В частности, число вершин нечётной степени чётно.
Для орграфа принимаем знаки со с. 14 лекций:
\(d^+(v)=|\{u:uv\in E\}|\) — число \emph{входящих},
\(d^-(v)=|\{u:vu\in E\}|\) — число \emph{исходящих} дуг.
В других источниках знаки часто противоположны; здесь везде используется
указанное лекционное соглашение.
\(\sum d^+(v)=\sum d^-(v)=|E|\), поскольку каждая дуга
имеет ровно одно начало и один конец.

\textbf{Линейная алгебра.} Используются определитель, его разложение по строке,
мультилинейность и правило умножения матриц. Необходимая для билета 7
формула Бине–Коши обосновывается непосредственно там.
Для плоских графов используется теорема Жордана: простая замкнутая кривая
разделяет плоскость на внутреннюю и внешнюю области и служит их общей границей.
Это топологический исходный факт, а не дополнительная графовая лемма.

\textbf{Обозначения матриц и раскрасок.}
Матрица смежности — \(A(G)\), матрица Кирхгофа — \(K(G)\).
Хроматический многочлен обозначается \(\chi_G(t)\), как в лекциях;
его аргумент может обозначаться также \(k\). Хроматическое число —
\(\chi(G)\), без нижнего индекса: это число, а не многочлен.

\endgroup
\ticket{Матрица смежности и её степени}
\label{ticket:1}
\begin{definition}[Матрица смежности]
Для графа с фиксированным порядком вершин \(v_1,\ldots,v_n\) матрица
смежности \(A=(a_{ij})\) задаётся равенством \(a_{ij}=1\), если \(v_i\)
и \(v_j\) смежны, и \(a_{ij}=0\) иначе. У орграфа \(a_{ij}=1\)
означает наличие дуги \(v_i\to v_j\).
Для графа с кратными рёбрами \(a_{ij}\) равно их числу; для орграфа
с петлями \(a_{ii}\) равно числу петель. Далее неориентированные графы
с петлями в матричных формулах не рассматриваются.
\end{definition}
У неориентированного графа \(A=A^{\mathsf T}\); у простого графа
\(a_{ii}=0\), \(d(v_i)=\sum_j a_{ij}\). В орграфе суммы строк дают
исходящие, а суммы столбцов — входящие степени.
\begin{graphillustration}{Пример: граф $P_3$ и его матрица в порядке $v_1,v_2,v_3$.}{fig:01-adjacency}
\begin{tikzpicture}[font=\small]
\node[gi vertex] (v1) at (0,0) {$v_1$};
\node[gi vertex] (v2) at (1.5,1) {$v_2$};
\node[gi vertex] (v3) at (3,0) {$v_3$};
\draw[gi edge] (v1) -- (v2);
\draw[gi edge] (v2) -- (v3);
\node[gi note] at (6,0) {$A=\begin{pmatrix}0&1&0\\1&0&1\\0&1&0\end{pmatrix}$};
\node[gi note] at (6,-1.25) {Строка $i$ — соседи $v_i$};
\end{tikzpicture}
\end{graphillustration}


\begin{lemma}[Разбиение по предпоследней вершине]
Пусть \(w_k(i,j)\) — число маршрутов длины \(k\) из \(v_i\) в \(v_j\).
Тогда
\[
 w_0(i,j)=\begin{cases}1,&i=j,\\0,&i\ne j,\end{cases}
 \qquad w_{k+1}(i,j)=\sum_{s=1}^n w_k(i,s)a_{sj}.
\]
\end{lemma}
\begin{proof}
Маршрут нулевой длины состоит из одной вершины. У каждого маршрута
длины \(k+1\) однозначно определены первые \(k\) рёбер и последнее ребро.
Если предпоследняя вершина \(v_s\), выборов этих частей соответственно
\(w_k(i,s)\) и \(a_{sj}\). Разные значения \(s\) дают непересекающиеся классы.
\end{proof}
\begin{theorem}[Комбинаторный смысл степеней матрицы]
Для любого целого \(k\ge0\) элемент \((A^k)_{ij}\) равен числу маршрутов
длины \(k\) из \(v_i\) в \(v_j\).
\end{theorem}
\begin{proof}
Индукция по \(k\). При \(k=0\) имеем \(A^0=I\), что совпадает
с \(w_0(i,j)\). Если утверждение верно для \(k\), то
\[
 (A^{k+1})_{ij}=\sum_s(A^k)_{is}a_{sj}
 =\sum_s w_k(i,s)a_{sj}=w_{k+1}(i,j)
\]
по лемме. Индукционный переход доказан.
\end{proof}
\begin{cor}
Вершина \(v_j\) достижима из \(v_i\) тогда и только тогда, когда
\((I+A+\cdots+A^{n-1})_{ij}>0\).
\end{cor}
\begin{proof}
Положительное слагаемое даёт маршрут, значит, достижимость.
Обратно, из маршрута удалением повторений получается простой путь;
при различных концах в нём не более \(n-1\) рёбер. При совпадающих
концах достаточно маршрута длины ноль.
\end{proof}
\begin{example}[Почему это не число простых путей]
В графе из одного ребра \(uv\) существует маршрут \(u,v,u\) длины два,
поэтому \((A^2)_{uu}=1\). Простого пути длины два из \(u\) в \(u\) нет.
В лекциях такой маршрут называется путём; простота пути оговаривается отдельно.
\end{example}

\ticket{Подсчёт маршрутов с помощью матрицы смежности}
\label{ticket:2}
\begin{definition}[Матричный подсчёт маршрутов]
Обозначим \(w_k(i,j)=(A^k)_{ij}\). В отличие от простого пути, маршрут
допускает многократное посещение вершин и повторное прохождение рёбер.
Для кратных рёбер различается также выбор конкретного ребра.
\end{definition}
\begin{lemma}[Разбиение маршрута на две части]
Для \(p,q\ge0\)
\[
 w_{p+q}(i,j)=\sum_s w_p(i,s)w_q(s,j).
\]
\end{lemma}
\begin{proof}
Разрежем маршрут после \(p\)-го ребра. Вершина разреза \(v_s\) и обе
части определены однозначно. Обратно, любую пару таких частей можно
склеить. Для фиксированного \(s\) действует правило произведения,
затем правило суммы. Это комбинаторное доказательство \(A^{p+q}=A^pA^q\).
\end{proof}
\begin{graphillustration}{Каждый маршрут однозначно разбивается в $v_s$ на префикс длины $p$ и суффикс длины $q$; линии обозначают целые участки.}{fig:02-split}
\begin{tikzpicture}[font=\small]
\node[gi vertex] (i) at (0,0) {$v_i$};
\node[gi selected] (s) at (4,0) {$v_s$};
\node[gi vertex] (j) at (8,0) {$v_j$};
\draw[gi emphasis] (i) to[bend left=15] (s);
\draw[gi alternate] (s) to[bend left=15] (j);
\node[gi note] at (2,0.75) {$p$ рёбер};
\node[gi note] at (6,0.75) {$q$ рёбер};
\node[gi note] at (4,-0.8) {разрез после $p$-го ребра};
\end{tikzpicture}
\end{graphillustration}


\begin{theorem}[Полезные суммы]
Число маршрутов длины \(k\), начинающихся в \(v_i\), равно
\(\sum_j(A^k)_{ij}\); число всех маршрутов длины \(k\) равно
\(\boldsymbol1^{\mathsf T}A^k\boldsymbol1\); число замкнутых маршрутов
длины \(k\) с отмеченной начальной вершиной равно \(\operatorname{tr}A^k\).
\end{theorem}
\begin{proof}
В первом случае маршруты разбиты по последней вершине, во втором —
по упорядоченной паре концов. В третьем берутся только маршруты,
у которых концы совпадают, то есть диагональные элементы матрицы.
Основная теорема билета 1 превращает каждое слагаемое в соответствующее число маршрутов.
\end{proof}
\begin{cor}[Для простого неориентированного графа]
\[
 (A^2)_{ii}=d(v_i),\qquad (A^2)_{ij}=|N(v_i)\cap N(v_j)|\quad(i\ne j),
 \qquad \operatorname{tr}A^2=2m.
\]
Число треугольников равно \(\operatorname{tr}A^3/6\).
\end{cor}
\begin{proof}
Двухшаговый маршрут имеет вид \(v_i,v_s,v_j\); его середина — общий
сосед концов. При совпадении концов можно выбрать любого соседа.
Сумма степеней равна \(2m\). Замкнутый трёхшаговый маршрут в простом
графе проходит по трём различным вершинам треугольника; у каждого
треугольника три выбора начала и два направления обхода.
\end{proof}
\begin{example}[Треугольник]
Для \(K_3\)
\[
 A=\begin{pmatrix}0&1&1\\1&0&1\\1&1&0\end{pmatrix},\quad
 A^2=\begin{pmatrix}2&1&1\\1&2&1\\1&1&2\end{pmatrix},\quad
 A^3=\begin{pmatrix}2&3&3\\3&2&3\\3&3&2\end{pmatrix}.
\]
Из каждой вершины выходят четыре маршрута длины два; всего их 12.
След \(A^3\) равен 6, так что треугольник ровно один.
\end{example}
\begin{graphillustration}{Шесть замкнутых маршрутов длины 3 одного треугольника: три начала (двойной контур) и два направления.}{fig:02-triangle}
\begin{tikzpicture}[font=\small]
\begin{scope}[shift={(0.0,0.0)}]
\node[gi selected] (1) at (0,0) {$1$};
\node[gi vertex] (2) at (1,1.1) {$2$};
\node[gi vertex] (3) at (2,0) {$3$};
\draw[gi arrow] (1) -- (2);
\draw[gi arrow] (2) -- (3);
\draw[gi arrow] (3) -- (1);
\node[gi note] at (1,-0.65) {$1231$};
\end{scope}
\begin{scope}[shift={(4.2,0.0)}]
\node[gi vertex] (1) at (0,0) {$1$};
\node[gi selected] (2) at (1,1.1) {$2$};
\node[gi vertex] (3) at (2,0) {$3$};
\draw[gi arrow] (2) -- (3);
\draw[gi arrow] (3) -- (1);
\draw[gi arrow] (1) -- (2);
\node[gi note] at (1,-0.65) {$2312$};
\end{scope}
\begin{scope}[shift={(8.4,0.0)}]
\node[gi vertex] (1) at (0,0) {$1$};
\node[gi vertex] (2) at (1,1.1) {$2$};
\node[gi selected] (3) at (2,0) {$3$};
\draw[gi arrow] (3) -- (1);
\draw[gi arrow] (1) -- (2);
\draw[gi arrow] (2) -- (3);
\node[gi note] at (1,-0.65) {$3123$};
\end{scope}
\begin{scope}[shift={(0.0,-2.8)}]
\node[gi selected] (1) at (0,0) {$1$};
\node[gi vertex] (2) at (1,1.1) {$2$};
\node[gi vertex] (3) at (2,0) {$3$};
\draw[gi arrow] (1) -- (3);
\draw[gi arrow] (3) -- (2);
\draw[gi arrow] (2) -- (1);
\node[gi note] at (1,-0.65) {$1321$};
\end{scope}
\begin{scope}[shift={(4.2,-2.8)}]
\node[gi vertex] (1) at (0,0) {$1$};
\node[gi selected] (2) at (1,1.1) {$2$};
\node[gi vertex] (3) at (2,0) {$3$};
\draw[gi arrow] (2) -- (1);
\draw[gi arrow] (1) -- (3);
\draw[gi arrow] (3) -- (2);
\node[gi note] at (1,-0.65) {$2132$};
\end{scope}
\begin{scope}[shift={(8.4,-2.8)}]
\node[gi vertex] (1) at (0,0) {$1$};
\node[gi vertex] (2) at (1,1.1) {$2$};
\node[gi selected] (3) at (2,0) {$3$};
\draw[gi arrow] (3) -- (2);
\draw[gi arrow] (2) -- (1);
\draw[gi arrow] (1) -- (3);
\node[gi note] at (1,-0.65) {$3213$};
\end{scope}
\end{tikzpicture}
\end{graphillustration}


\textbf{Граница применимости.} Деление следа на \(2k\) при \(k\ge4\)
обычно не даёт число простых циклов: появляются замкнутые маршруты
с повторениями. Для орграфа формула для треугольников также меняется.

\ticket{Сильная связность, порядок и конденсация}
\label{ticket:3}
\begin{definition}[Сильная связность]
Пишем \(u\leadsto v\), если существует направленный маршрут из \(u\)
в \(v\), включая маршрут длины ноль. Вершины сильно связаны,
\(u\sim v\), если \(u\leadsto v\) и \(v\leadsto u\).
Орграф сильно связен, если все его вершины попарно сильно связаны.
\end{definition}
\begin{theorem}
Сильная связность — отношение эквивалентности. Его классы называются
компонентами сильной связности (КСС); каждый класс индуцирует сильно связный подграф.
\end{theorem}
\begin{proof}
Рефлексивность следует из нулевого маршрута; симметричность заложена
в определении. Если \(u\sim v\) и \(v\sim w\), склейка маршрутов даёт
\(u\leadsto w\) и \(w\leadsto u\), то есть транзитивность.
Если \(u,v\) лежат в одной КСС, то каждая вершина \(x\) маршрута
\(u\leadsto v\) достижима из \(u\) и достигает \(u\) через \(v\).
Значит, \(x\) принадлежит той же КСС и маршрут не выходит из класса.
\end{proof}
\begin{definition}[Слабая связность]
Орграф слабо связен, если связен его неориентированный граф, полученный
забыванием направлений. Сильная связность влечёт слабую, обратное неверно:
одна дуга на двух вершинах даёт контрпример.
Ациклический орграф не имеет направленных циклов положительной длины.
\end{definition}
\begin{theorem}[Топологическая сортировка]
Орграф ацикличен тогда и только тогда, когда существует биекция
\(\varphi:V\to\{1,\ldots,n\}\), для которой
\(uv\in E\Rightarrow\varphi(u)<\varphi(v)\).
\end{theorem}
\begin{proof}
В непустом ациклическом орграфе есть вершина без входящих дуг.
Иначе, выбирая предшественника, затем его предшественника и так далее,
в конечном множестве получим повторение вершины и направленный цикл.
Ставим вершину без входящих дуг первой, удаляем её и повторяем в оставшемся
ациклическом орграфе. При удалении вершины все её предшественники уже
стоят в списке, значит, любая дуга идёт вперёд.
Обратно, вдоль направленного цикла значения \(\varphi\) должны были бы
строго возрастать и вернуться к исходному значению, что невозможно.
\end{proof}
\begin{cor}[Достижимость как порядок]
В ациклическом орграфе отношение \(u\preceq v\), заданное существованием
пути из \(u\) в \(v\), включая нулевой, является частичным порядком.
\end{cor}
\begin{proof}
Рефлексивность даёт нулевой путь, транзитивность — склейка путей.
Взаимная достижимость разных вершин давала бы направленный цикл,
поэтому отношение антисимметрично.
\end{proof}
\textbf{Дополнение по списку вопросов.} В скане следующие определения
не выписаны; они дополняют лекционный порядок на ациклическом орграфе.
\begin{definition}[Конденсация и порядок]
Конденсация \(D^*\) имеет вершинами КСС; дуга \(C\to B\), \(C\ne B\),
проводится, если в исходном орграфе есть дуга из \(C\) в \(B\).
На КСС зададим \(C\preceq B\), если \(C=B\) либо в конденсации есть
направленный путь из \(C\) в \(B\).
\end{definition}
\begin{theorem}
Конденсация ациклична, а \(\preceq\) является частичным порядком.
\end{theorem}
\begin{proof}
Если различные КСС образуют ориентированный цикл, то, используя
сильную связность внутри классов и дуги между ними, из любой вершины
их объединения можно попасть в любую другую. Все они должны составлять
один класс — противоречие. Рефлексивность и транзитивность порядка
следуют из достижимости. Если \(C\preceq B\preceq C\), то при \(C\ne B\)
получается замкнутый направленный маршрут и, после удаления повторений,
цикл в конденсации. Значит, \(C=B\), что даёт антисимметричность.
\end{proof}
\begin{definition}[Диаграмма порядка]
Диаграмма Хассе содержит дугу \(C\to B\), если \(C\prec B\), но нет
класса \(F\), для которого \(C\prec F\prec B\).
Это отношение покрытия, а не все сравнимые пары.
\end{definition}
\begin{lemma}
Диаграмма Хассе ациклична и имеет ту же достижимость, что и конденсация.
\end{lemma}
\begin{proof}
Цикл противоречит строгому возрастанию в частичном порядке.
Для \(C\prec B\) выберем цепочку между ними максимальной длины;
конечность множества гарантирует её существование. Соседние элементы
цепочки образуют покрытие, иначе можно вставить промежуточный элемент.
Поэтому сравнимость реализуется путём по диаграмме. Обратное следует из транзитивности.
\end{proof}
\begin{center}
\begin{tikzpicture}[>=Stealth,every node/.style={circle,draw,minimum size=6mm}]
\node (a) at (0,0) {$A$};\node (b) at (2,0) {$B$};\node (c) at (4,0) {$C$};
\draw[->] (a)--(b);\draw[->] (b)--(c);\draw[->,bend left=35] (a) to (c);
\end{tikzpicture}
\end{center}
Здесь конденсация может содержать все три дуги; в диаграмме Хассе
остаётся только \(A\to B\to C\). Конденсация и диаграмма Хассе —
разные графы; оба ацикличны, но первая может содержать транзитивные дуги.

\ticket{Рёберная двусвязность}
\label{ticket:4}
\begin{definition}[Рёберная двусвязность]
Вершины \(u,v\) рёберно двусвязны, \(u\sim_E v\), если существуют
два \(uv\)-пути без общих рёбер. Для \(u=v\) допускаются два нулевых пути.
Пути можно считать простыми, удалив замкнутые участки.
\end{definition}
\begin{theorem}
Рёберная двусвязность — отношение эквивалентности на вершинах.
Его классы называются компонентами рёберной двусвязности.
\end{theorem}
\begin{proof}
Рефлексивность следует из нулевых путей, симметричность — из обращения.
Пусть \(u\sim_E v\), \(v\sim_E w\). Возьмём непересекающиеся по рёбрам
простые пути \(P_1,P_2\) из \(u\) в \(v\) и такие же пути \(Q_1,Q_2\)
из \(w\) в \(v\). Пусть \(a,b\) — первые вершины встречи \(Q_1,Q_2\)
с \(P_1\cup P_2\). До них пути не используют рёбер \(P_1,P_2\).
Если \(w\) уже лежит на объединении, допускаются нулевые начальные участки.

Если \(a,b\) можно отнести к разным \(P_i\), склеим начальный участок
каждого \(Q_i\) с участком соответствующего \(P_i\) до \(u\).
Получены два пути без общих рёбер. Иначе обе точки лежат только
на одном пути, пусть \(P_1\), в порядке \(u,\ldots,a,\ldots,b,\ldots,v\)
(при необходимости поменяем \(Q_1,Q_2\) местами).
Первый путь идёт из \(w\) по \(Q_1\) до \(a\), затем по \(P_1\) к \(u\).
Второй — по \(Q_2\) до \(b\), по \(P_1\) от \(b\) к \(v\), затем
по \(P_2\) к \(u\). Участки \(P_1\) не имеют общих рёбер, как и
начальные участки \(Q_1,Q_2\). При \(a=b\) рассуждение также работает.
Внутри отдельного пути возможные повторения удаляются.
Таким образом, \(u\sim_E w\).
\end{proof}
\begin{graphillustration}{Два случая в доказательстве транзитивности теоремы 4.1. Жирная сплошная и штриховая линии выделяют полученные $uw$-пути $R_1,R_2$ без общих рёбер; тонкие серые линии — неиспользуемые участки. Показаны только начальные участки $Q_1,Q_2$ до первых встреч с $P_1\cup P_2$.}{fig:04-transitivity}
\begin{tikzpicture}[font=\small]
\begin{scope}
  \node[gi note] at (2,2) {1. Точки встречи на разных путях};
  \node[gi vertex] (u) at (0,0) {$u$};
  \node[gi vertex] (v) at (4,0) {$v$};
  \node[gi vertex] (a) at (2,1.1) {$a$};
  \node[gi vertex] (b) at (2,-1.1) {$b$};
  \node[gi vertex] (w) at (2,0) {$w$};
  \draw[gi edge,black!45] (a)--(v);
  \draw[gi edge,black!45] (b)--(v);
  \draw[gi emphasis] (w)--(a)--(u);
  \draw[gi alternate] (w)--(b)--(u);
  \node[gi note] at (3.25,.9) {$P_1$};
  \node[gi note] at (3.25,-.9) {$P_2$};
  \node[gi note] at (2.45,.55) {$Q_1$};
  \node[gi note] at (2.45,-.55) {$Q_2$};
  \node[gi note] at (2,-2.9) {$R_1:w\xrightarrow{Q_1}a\xrightarrow{P_1}u$};
  \node[gi note] at (2,-3.55) {$R_2:w\xrightarrow{Q_2}b\xrightarrow{P_2}u$};
\end{scope}
\begin{scope}[shift={(7.2,0)}]
  \node[gi note] at (2,2) {2. Обе точки встречи на $P_1$};
  \node[gi vertex] (u) at (0,0) {$u$};
  \node[gi vertex] (v) at (4,0) {$v$};
  \node[gi vertex] (a) at (1.3,-.8) {$a$};
  \node[gi vertex] (b) at (2.7,-.8) {$b$};
  \node[gi vertex] (w) at (2,-2.1) {$w$};
  \draw[gi edge,black!45] (a)--(b);
  \draw[gi emphasis] (w)--(a)--(u);
  \draw[gi alternate] (w)--(b)--(v);
  \draw[gi alternate] (v) to[bend right=35] (u);
  \node[gi note] at (2,1.35) {$P_2$};
  \node[gi note] at (2,-.4) {$P_1$};
  \node[gi note] at (1.15,-1.6) {$Q_1$};
  \node[gi note] at (2.85,-1.6) {$Q_2$};
  \node[gi note] at (2,-2.9) {$R_1:w\xrightarrow{Q_1}a\xrightarrow{P_1}u$};
  \node[gi note] at (2,-3.55) {$R_2:w\xrightarrow{Q_2}b\xrightarrow{P_1}v\xrightarrow{P_2}u$};
\end{scope}
\end{tikzpicture}
\end{graphillustration}


\begin{definition}[Мост]
Ребро \(ab\) называется мостом, если его концы лежат в разных компонентах
рёберной двусвязности.
\end{definition}
\begin{lemma}[Мост, цикл и удаление]\label{lem:bridge}
Для ребра \(e=ab\) эквивалентны: оно является мостом; не лежит
ни на одном простом цикле; его удаление увеличивает число компонент связности.
\end{lemma}
\begin{proof}
Если \(e\) лежит на простом цикле, ребро и остальная дуга цикла дают
два пути между концами без общих рёбер, и удаление \(e\) не нарушает
связности: любое его использование заменяется обходом по циклу.
Если \(a\sim_E b\), хотя бы один из двух простых путей не использует
\(e\); вместе с \(e\) он образует простой цикл.
Наконец, если \(G-e\) всё ещё соединяет \(a,b\), простой путь между
ними вместе с \(e\) образует цикл. Если не соединяет, исходная компонента
распадается на две: любой простой путь от её вершины к \(a\) либо не использует
\(e\), либо до его прохождения соединяет вершину с \(b\).
\end{proof}
\begin{graphillustration}{Пример: удаление моста разделяет два треугольника. Пунктир отмечает удалённое ребро.}{fig:04-bridge}
\begin{tikzpicture}[font=\small]
\node[gi vertex] (a) at (0,0) {$a$};
\node[gi vertex] (p) at (-1,0.8) {$p$};
\node[gi vertex] (q) at (-1,-0.8) {$q$};
\node[gi vertex] (b) at (1.5,0) {$b$};
\node[gi vertex] (r) at (2.5,0.8) {$r$};
\node[gi vertex] (s) at (2.5,-0.8) {$s$};
\draw[gi edge] (a) -- (p);
\draw[gi edge] (p) -- (q);
\draw[gi edge] (q) -- (a);
\draw[gi edge] (b) -- (r);
\draw[gi edge] (r) -- (s);
\draw[gi edge] (s) -- (b);
\draw[gi emphasis] (a) -- (b);
\node[gi note] at (0.75,1.25) {$e=ab$ — мост};
\node[gi note] at (0.75,-1.45) {$c(G)=1$};
\node[gi note] at (4.3,0) {$\longrightarrow$};
\begin{scope}[shift={(6.5,0)}]
\node[gi vertex] (a) at (0,0) {$a$};
\node[gi vertex] (p) at (-1,0.8) {$p$};
\node[gi vertex] (q) at (-1,-0.8) {$q$};
\node[gi vertex] (b) at (1.5,0) {$b$};
\node[gi vertex] (r) at (2.5,0.8) {$r$};
\node[gi vertex] (s) at (2.5,-0.8) {$s$};
\draw[gi edge] (a) -- (p);
\draw[gi edge] (p) -- (q);
\draw[gi edge] (q) -- (a);
\draw[gi edge] (b) -- (r);
\draw[gi edge] (r) -- (s);
\draw[gi edge] (s) -- (b);
\draw[gi removed] (a) -- (b);
\node[gi note] at (0.75,-1.45) {$c(G-e)=2$};
\end{scope}
\end{tikzpicture}
\end{graphillustration}


\begin{theorem}[Удаление одного ребра и компоненты]
Вершины рёберно двусвязны тогда и только тогда, когда связаны и остаются
связанными после удаления любого одного ребра. Классы \(\sim_E\)
совпадают с компонентами графа после удаления всех мостов.
\end{theorem}
\begin{proof}
Два непересекающихся по рёбрам пути не разрушаются удалением одного ребра.
Если простой путь между \(u,v\) не проходит через мосты, концы каждого
его ребра рёберно двусвязны; по транзитивности \(u\sim_E v\).
Если же \(u,v\) находятся в разных компонентах после удаления мостов,
простой путь между ними проходит через некоторый мост.
Удаление этого моста разрывает связь между концами пути:
иначе его две стороны соединялись бы обходом. Это доказывает оба утверждения.
\end{proof}

\ticket{Вершинная двусвязность}
\label{ticket:5}
\source{Лекции, с. 7–9, 57; элементарные детали доказательств восстановлены составителем}
\begin{definition}[Вершинная двусвязность рёбер]
Два различных ребра \(ab,cd\) вершинно двусвязны, если их концы можно
попарно соединить двумя путями без общих вершин, не используя сами
эти рёбра. Разрешены оба сопоставления концов; для смежных рёбер один
путь может быть нулевым. Каждое ребро считаем двусвязным с самим собой.
Эквивалентно: два различных ребра лежат на общем простом цикле.
Действительно, удаление двух рёбер из цикла даёт нужные пути, а их
склейка с этими рёбрами даёт цикл.
\end{definition}
\textbf{Элементарные детали, используемые далее.}
Граф на \(n\ge3\) вершинах назовём двусвязным, если он связен
и удаление любой одной вершины оставляет его связным.
Следующая лемма раскрывает рисунки о непересекающихся путях без обращения к потокам.
\begin{lemma}[Два выхода к циклу]\label{lem:vertex-paths}
В двусвязном графе из вершины \(x\) к множеству \(S\), \(|S|\ge2\),
\(x\notin S\), идут два пути с разными концами в \(S\), пересекающиеся
только в \(x\) и не встречающие \(S\) внутренними вершинами.
Любые две вершины двусвязного графа лежат на общем простом цикле.
\end{lemma}
\begin{proof}
В связном графе сначала построим дерево из пути от \(x\) к \(S\)
и пути от другой вершины \(S\) к первому пути. После обрезания лишнего
это три ветви от вершины \(z\): к \(x\), \(s_1\), \(s_2\), причём
\(s_1,s_2\in S\), \(s_1\ne s_2\). Одна ветвь может быть нулевой.
Ветви обрезаем при первой встрече с \(S\), так что внутри ветви \(xz\)
нет вершин \(S\). Среди таких деревьев минимизируем длину ветви \(xz\).
Если \(z\ne x\), в \(G-z\) есть путь от \(x\) к одной из двух остальных
ветвей, так как хотя бы одна из \(s_1,s_2\) отлична от \(z\).
Берём его участок от последней встречи \(t\) с \(xz\) до первой
встречи с другой ветвью. Новый обход, ветвящийся в \(t\), даёт две
разные конечные вершины \(S\) и более короткую общую ветвь \(xt\).
Противоречие. Значит, \(z=x\); получены требуемые пути.

Для второго утверждения сначала построим цикл через \(u\): у \(u\)
есть два разных соседа, соединённых простым путём в \(G-u\).
Если \(v\) вне этого цикла, применим первую часть с \(x=v\), \(S=V(C)\).
Два выхода к \(a,b\in C\) вместе с той дугой \(C\), которая содержит
\(u\), образуют простой цикл через \(u,v\).
\end{proof}
\begin{lemma}[Циклы через рёбра]\label{lem:blocks}
В двусвязном графе любые два различных ребра лежат на общем простом цикле.
Объединение двух двусвязных подграфов с двумя общими вершинами двусвязно.
\end{lemma}
\begin{proof}
Последнее утверждение следует из того, что после удаления одной вершины
оба подграфа связны и у них остаётся общая вершина.
Для первого подразделим два выбранных ребра новыми вершинами \(x,y\).
Подразделение сохраняет двусвязность: удаление старой вершины оставляет
новую прикреплённой к другому концу ребра, а удаление новой вершины
равносильно удалению немостового ребра. В двусвязном графе мостов нет:
при \(n\ge3\) хотя бы один конец моста был бы разделяющей вершиной.
По предыдущей лемме \(x,y\) лежат на общем цикле. Убирая подразделения,
получаем цикл через оба исходных ребра.
\end{proof}
\begin{theorem}[Отношение и его классы]
Вершинная двусвязность является отношением эквивалентности на \emph{рёбрах}.
Её классы называются компонентами вершинной двусвязности, или блоками;
сам блок рассматривается вместе со всеми концами его рёбер.
\end{theorem}
\begin{proof}
Рефлексивность задана, симметричность очевидна. Пусть различные пары
\(e,f\) и \(f,g\) лежат на циклах \(C,D\). Циклы двусвязны, у них общими
являются оба конца \(f\), поэтому \(C\cup D\) двусвязен.
По предыдущей лемме \(e,g\) лежат на общем цикле. Если какие-то рёбра
совпали, транзитивность очевидна. Следовательно, классы разбивают \(E\).
Немостовые классы — максимальные двусвязные подграфы: объединяем
циклы через фиксированное ребро класса, применяя лемму об объединении;
обратно, все рёбра двусвязного подграфа принадлежат одному классу.
Мост образует одноэлементный класс.
\end{proof}
Для удобства разложения графа будем также называть изолированную вершину
вырожденным блоком. Это дополнительное соглашение: класса рёбер у неё нет.
Блоки-мосты также называем вырожденными; остальные блоки двусвязны
и называются невырожденными.
Разные блоки имеют не более одной общей вершины, иначе их объединение
двусвязно и классы совпадают. Каждое ребро принадлежит ровно одному блоку.
\begin{definition}[Точка сочленения]
Вершина, принадлежащая нескольким блокам с рёбрами, называется точкой сочленения.
\end{definition}
\begin{theorem}[Критерий точки сочленения]
Вершина \(v\) является точкой сочленения тогда и только тогда,
когда её удаление увеличивает число компонент связности.
\end{theorem}
\begin{proof}
Достаточно рассмотреть компоненту, содержащую \(v\).
Если два инцидентных \(v\) ребра \(va,vb\) принадлежат разным блокам,
в \(G-v\) нет пути от \(a\) к \(b\): простой такой путь вместе
с \(av,vb\) дал бы общий цикл. Поэтому удаление разъединяет компоненту.
Обратно, если она распалась, у \(v\) есть соседи \(a,b\) в разных
полученных компонентах. Общий цикл через \(va,vb\) дал бы обходной
путь от \(a\) к \(b\) без \(v\), что невозможно. Значит, блоки разные.
Изолированная вершина и конец висячего ребра проверяются непосредственно.
\end{proof}
\begin{cor}[Критерий двусвязности]
При \(n\ge3\) эквивалентны: связность без точек сочленения;
связность после удаления любой вершины; наличие общего простого цикла
для любых двух различных вершин.
\end{cor}
\begin{proof}
Первую эквивалентность даёт критерий удаления. Вторая следует из леммы
о двух выходах; в обратную сторону цикл сохраняет связь двух оставшихся
вершин после удаления любой третьей.
\end{proof}
На вершинах отношение «лежать в одном блоке» не транзитивно:
два треугольника с общей вершиной дают контрпример. Это согласуется
с лекционным определением отношения именно на рёбрах.
\begin{lemma}[Дерево блоков]\label{lem:block-tree}
Для связного графа двудольный граф инцидентности блоков и точек
сочленения является деревом. В частности, если блоков больше одного,
есть не менее двух крайних блоков, каждый из которых содержит ровно
одну точку сочленения исходного графа.
\end{lemma}
\begin{proof}
Сначала построим граф инцидентности \(T\) всех блоков и \emph{всех}
вершин исходного графа. Ребро в \(T\) означает принадлежность вершины блоку.
Он связен: любое ребро исходного графа лежит в блоке, поэтому путь
исходного графа превращается в чередующийся маршрут по \(T\).

Пусть в \(T\) есть простой цикл. Рассмотрим объединение его блоков;
в каждом отмечены две различные вершины стыка с соседними блоками.
После удаления любой одной вершины \(x\) каждый невырожденный блок
остаётся связным; блок-мост либо сохраняется, либо от него остаётся
одна вершина стыка. Блок из изолированной вершины в цикл входить не может.
Среди различных отмеченных стыков исчезает не более одного,
поэтому вдоль остальных стыков блоки всё ещё соединяются в цепочку
или цикл. Следовательно, всё объединение без \(x\) связно.
Это двусвязный подграф, содержащий несколько различных блоков,
вопреки их максимальности (а мост вообще не может лежать в двусвязном подграфе).
Значит, \(T\) — дерево.

Если вершина \(v\) принадлежит двум или более блокам, её узел в \(T\)
имеет степень не менее двух. Удаление этого узла разъединяет дерево.
Путь в \(G-v\) между вершинами разных ветвей дал бы маршрут в
\(T-v\), что невозможно. Значит, \(v\) — точка сочленения.
Если же \(v\) принадлежит единственному блоку, его удаление не нарушает
связности остального графа: невырожденный блок остаётся связным,
а в блоке-мосте такая вершина — лист исходного графа.
Одноточечный граф рассматривается непосредственно.

Удалим из \(T\) все листовые узлы обычных вершин; останется именно
граф инцидентности блоков и точек сочленения. Он является деревом
(при единственном блоке — одной вершиной).
Дерево с более чем одной вершиной имеет не менее двух листьев:
это концы максимального простого пути. Узел точки сочленения
имеет степень не менее двух, поэтому листья соответствуют блокам.
\end{proof}

\ticket{Леса и эквивалентные определения дерева}
\label{ticket:6}
\source{Лекции, с. 15–21; Карпов, §1.5.2, с. 25–27}
\begin{definition}[Лес, дерево, лист]
Ациклический неориентированный граф (без простых циклов) называется лесом.
Дерево — связный ациклический граф. Лист — вершина степени 1.
Остовное дерево — дерево, являющееся остовным подграфом исходного графа.
\end{definition}
\begin{lemma}[Листья]
В дереве с не менее чем двумя вершинами есть хотя бы два листа.
Удаление листа с его ребром оставляет дерево.
\end{lemma}
\begin{proof}
Возьмём простой путь максимальной длины. У каждого его конца нет соседа
вне пути (путь удлинился бы) и нет другого соседа на пути, кроме соседней
вершины (возник бы цикл). Оба конца — листья.
Путь между оставшимися вершинами не проходит через лист внутренней
вершиной, поэтому после удаления связность сохраняется; циклов не возникает.
\end{proof}
Для одноточечного дерева листьев нет; утверждения о листьях требуют \(n\ge2\).
\begin{theorem}[Любые два из трёх свойств]
Для непустого графа любые два из следующих свойств влекут третье:
\begin{enumerate}[nosep]
\item \(m=n-1\);
\item граф связен;
\item граф ацикличен.
\end{enumerate}
Следовательно, любая пара характеризует дерево.
\end{theorem}
\begin{proof}
\(2,3\Rightarrow1\). Индукция по \(n\): при \(n=1\) рёбер нет.
При \(n>1\) удаляем лист. Для оставшегося дерева \(m'=n'-1\),
поэтому \(m=m'+1=(n-1)-1+1=n-1\).

\(1,3\Rightarrow2\). Каждая из \(c\) компонент — дерево.
Если размеры компонент \(n_1,\ldots,n_c\), по уже доказанному
\(m=\sum_i(n_i-1)=n-c\). Сравнивая с \(m=n-1\), получаем \(c=1\).

\(1,2\Rightarrow3\). Снова индукция по \(n\). При \(n=1\) очевидно.
При \(n>1\) степени положительны, а их сумма \(2n-2<2n\),
поэтому есть вершина степени 1. Удаление этого листа сохраняет связность:
простой путь между оставшимися вершинами не мог использовать лист внутри.
Остаётся связный граф на \(n-1\) вершинах с \(n-2\) рёбрами; по индукции
он ацикличен. Добавление листа цикла не создаёт.
\end{proof}
\begin{cor}[Число рёбер леса]\label{lem:forest}
Для леса \(m=n-c\); для связного графа \(m\ge n-1\).
\end{cor}
\begin{proof}
Первое равенство уже получено суммированием по компонентам.
Если связный граф имеет цикл, удаляем из него ребро, сохраняя связность.
Повторяем до дерева. В нём \(n-1\) рёбер, поэтому исходно их не меньше.
\end{proof}
\begin{theorem}[Единственность пути]
Граф — дерево тогда и только тогда, когда между любыми двумя различными
вершинами существует ровно один простой путь.
\end{theorem}
\begin{proof}
В дереве связность даёт путь. У двух разных простых путей возьмём
первое расхождение и следующую встречу: участки между ними образуют цикл.
Обратно, условие даёт связность, а простой цикл дал бы два пути
между любыми двумя своими различными вершинами.
\end{proof}
\begin{theorem}[Минимальная связность и максимальная ацикличность]
Для непустого графа эквивалентны: быть деревом; быть связным и терять
связность при удалении любого ребра; быть ацикличным и получать цикл
при добавлении любого отсутствующего ребра.
\end{theorem}
\begin{proof}
В дереве каждое ребро — мост, а добавление ребра замыкает единственный
путь между его концами в цикл. Обратно, минимально связный граф не может
иметь цикл, поскольку его ребро удалялось бы без потери связности.
Максимально ацикличный граф не может иметь две компоненты: ребро между
ними цикла не создаёт. В обоих случаях получено дерево.
\end{proof}
\begin{cor}[Вершины и остов]
В дереве каждая вершина степени больше 1 является точкой сочленения.
Каждый связный граф имеет остовное дерево.
\end{cor}
\begin{proof}
После удаления нелистовой вершины её разные соседи не могут остаться
связанными: такой обход дал бы цикл в дереве. Существование остова
доказывается последовательным удалением рёбер циклов, как выше.
\end{proof}
\textbf{Уточнение записи с. 19.} Обратное утверждение о точках сочленения
не используется: треугольник с листом у каждой вершины имеет все нелистовые
вершины точками сочленения, но деревом не является.

\ticket{Матрица Кирхгофа и остовные деревья}
\label{ticket:7}
\source{Лекции, с. 21–27; Карпов, §9.2, с. 425–428; детали Бине–Коши восстановлены составителем}
\begin{definition}[Матрица Кирхгофа]
Для неориентированного графа без петель, в котором допускаются кратные рёбра,
лапласиан (матрица Кирхгофа) равен \(K=D-A\), где \(D=\diag(d(v_1),\ldots,d(v_n))\).
Иными словами, \(K_{ii}=d(v_i)\), \(K_{ij}=-a_{ij}\) при \(i\ne j\).
Через \(K^{(r)}\) обозначим матрицу после удаления строки и столбца \(r\),
через \(\tau(G)\) — число остовных деревьев.
\end{definition}
\begin{definition}[Матрицы инцидентности]
Неориентированная матрица инцидентности имеет строки-вершины и столбцы-рёбра:
\(I_{ve}=1\), если \(v\) — конец \(e\), иначе \(0\).
Ориентированная матрица меняет у конца дуги единицу на \(-1\),
а у начала оставляет \(+1\). В доказательстве обозначаем её \(B\),
чтобы не путать с единичной матрицей.
\end{definition}
\begin{lemma}[Матрица инцидентности]
Произвольно ориентируем рёбра и составим матрицу \(B\) размера \(n\times m\):
в столбце дуги \(u\to v\) стоят \(+1\) у \(u\), \(-1\) у \(v\), прочие элементы нулевые.
Тогда \(BB^{\mathsf T}=K\). Если удалить строку \(r\), полученная матрица \(B_r\)
удовлетворяет \(B_rB_r^{\mathsf T}=K^{(r)}\).
\end{lemma}
\begin{proof}
Скалярный квадрат строки равен числу инцидентных рёбер — степени вершины.
Скалярное произведение разных строк получает вклад \(-1\) ровно от рёбер
между соответствующими вершинами. Это и есть элементы \(K\).
Удаление строки в обоих множителях даёт нужную главную подматрицу.
\end{proof}
\begin{lemma}[Миноры инцидентности]
Пусть \(S\) состоит из \(n-1\) рёбер и \(B_{r,S}\) — соответствующая
квадратная подматрица \(B_r\). Тогда
\[
 \det B_{r,S}=\begin{cases}\pm1,&(V,S)\text{ — дерево},\\0,&\text{иначе}.\end{cases}
\]
\end{lemma}
\begin{proof}
Если \((V,S)\) не дерево, при \(|S|=n-1\) он несвязен по билету 6.
Выберем компоненту, не содержащую \(r\). Сумма строк её вершин
в \(B_{r,S}\) равна нулю: у каждого выбранного ребра оба конца в этой
компоненте либо оба вне её. Строки линейно зависимы, определитель нулевой.
Для дерева доказываем модуль определителя индукцией: при \(n=1\)
определитель пустой матрицы равен 1. При \(n>1\) есть лист \(x\ne r\).
В его строке единственный ненулевой элемент \(\pm1\). Разложение по этой
строке сводит определитель к меньшему дереву с удалённым листом.
\end{proof}
\begin{lemma}[Бине–Коши в нужном виде]
Для матрицы \(M\) размера \(p\times q\), \(p\le q\),
\[
 \det(MM^{\mathsf T})=\sum_{S\subseteq[1,q],\ |S|=p}(\det M_{[:,S]})^2.
\]
\end{lemma}
\begin{proof}
Более общо, для \(X\) размера \(p\times q\) и \(Y\) размера \(q\times p\)
разложим каждый столбец \(XY\) как линейную комбинацию столбцов \(X\).
Мультилинейность определителя даёт сумму по \(p\)-кам индексов.
Повторяющиеся индексы дают нуль из-за равных столбцов. Для фиксированного
набора различных индексов \(S\) сумма по всем их перестановкам равна
\(\det X_{[:,S]}\det Y_{[S,:]}\): знаки перестановок возникают при
перестановке столбцов \(X\), а сумма произведений элементов \(Y\)
со знаками есть её определитель. При \(Y=X^{\mathsf T}\) оба определителя равны.
\end{proof}
\begin{theorem}[Кирхгоф]
Для любого \(r\) и \(n\ge1\)
\[
 \boxed{\tau(G)=\det K^{(r)}}.
\]
\end{theorem}
\begin{proof}
Если \(m<n-1\), остовных деревьев нет; ранг \(B_rB_r^{\mathsf T}\)
не больше \(m<n-1\), поэтому определитель нулевой.
Иначе по трём леммам
\[
 \det K^{(r)}=\det(B_rB_r^{\mathsf T})
 =\sum_{|S|=n-1}(\det B_{r,S})^2=\sum_{(V,S)\text{ — дерево}}1.
\]
Для \(n=1\) обе стороны равны 1. Доказательство также охватывает несвязные графы:
каждый минор в сумме нулевой и \(\tau(G)=0\).
\end{proof}
\begin{cor}[Все алгебраические дополнения]
При \(n\ge2\) для любых \(i,j\)
\( (-1)^{i+j}\det K_{\widehat i,\widehat j}=\tau(G)\).
\end{cor}
\begin{proof}
Если \(G\) связен, основной минор положителен, а \(K\boldsymbol1=0\).
Поэтому \(\rank K=n-1\) и ядро состоит из постоянных векторов.
Из \(K\operatorname{adj}K=0\) каждый столбец присоединённой матрицы
постоянен; её симметрия делает все элементы одинаковыми.
На диагонали они равны \(\tau(G)\). Если компонент не менее двух,
их индикаторы — независимые векторы ядра, так что \(\rank K\le n-2\)
и все миноры порядка \(n-1\) нулевые.
\end{proof}
\begin{example}[Лекционный граф на четырёх вершинах]
Возьмём рёбра \(12,23,34,41,13\): квадрат с одной диагональю.
Тогда
\[
K=\begin{pmatrix}3&-1&-1&-1\\-1&2&-1&0\\-1&-1&3&-1\\-1&0&-1&2\end{pmatrix},
\quad \tau(G)=\det\begin{pmatrix}2&-1&0\\-1&3&-1\\0&-1&2\end{pmatrix}=8.
\]
В предварительном вычислении на с. 21 рукописи у матрицы стоят обратные
знаки; здесь используется последующее определение \(K=D-A\), поэтому
определитель главного минора положителен.
\end{example}
\begin{cor}[Удаление–стягивание для числа остовов]
Для непетлевого ребра \(e\)
\[
\tau(G)=\tau(G-e)+\tau(G/e).
\]
\end{cor}
\begin{proof}
Остовы, не содержащие \(e\), — в точности остовы \(G-e\).
Стягивание \(e\) даёт биекцию между остовами, содержащими \(e\),
и остовами \(G/e\): обратная операция раскрывает стянутую вершину
и возвращает \(e\). Связность и отсутствие циклов при этом сохраняются.
После стягивания петли удаляются, но кратные рёбра сохраняются различимыми.
\end{proof}

\ticket{Коды Прюфера и формула Кэли}
\label{ticket:8}
\begin{definition}[Код Прюфера]
Вершины дерева помечены числами \(1,\ldots,n\), \(n\ge2\).
Пока остаётся больше двух вершин, удаляем лист с наименьшей меткой
и записываем метку его единственного соседа. Полученная последовательность
длины \(n-2\) называется кодом Прюфера.
\end{definition}
\begin{lemma}[Частоты меток]
В коде дерева метка \(v\) встречается ровно \(d_T(v)-1\) раз.
В каждом оставшемся дереве листья — ровно те его вершины,
которые отсутствуют в ещё не обработанном суффиксе кода.
\end{lemma}
\begin{proof}
Рёбра у \(v\) удаляются по одному. Пока \(v\) не удалена,
удаление соседнего листа записывает \(v\) и уменьшает её степень на 1.
Последнее инцидентное ребро исчезает при удалении самой \(v\)
или остаётся последним ребром дерева; в обоих случаях метка \(v\)
за него не записывается. Итого записей \(d_T(v)-1\).
То же рассуждение применимо к любому промежуточному дереву;
степень 1 равносильна нулю будущих записей.
\end{proof}
\textbf{Декодирование.} Для произвольного слова \(a_1,\ldots,a_{n-2}\)
поддерживаем множество не удалённых меток \(U\). На шаге \(t\)
выбираем наименьшую метку \(x\in U\), отсутствующую среди
\(a_t,\ldots,a_{n-2}\), проводим ребро \(xa_t\), удаляем \(x\) из \(U\)
и переходим к следующей позиции. Последние две вершины соединяем ребром.
\begin{lemma}[Корректность декодирования]
Алгоритм всегда выполним и строит дерево с заданным кодом.
\end{lemma}
\begin{proof}
Перед шагом \(t\) осталось \(n-t+1\) меток и \(n-t-1\) позиций кода.
Все метки суффикса ещё находятся в \(U\), поскольку ранее удалялись
только отсутствовавшие в более длинном суффиксе. Следовательно,
не менее двух меток \(U\) в суффиксе отсутствуют и нужная \(x\)
существует. При этом \(x\ne a_t\).

Если рассматривать шаги в обратном порядке, начинаем с последнего ребра
и каждый раз добавляем новую вершину одним ребром к уже построенному
графу: \(a_t\) в момент удаления \(x\) ещё не удалена.
Так сохраняются связность и отсутствие циклов, значит, построено дерево.
В оставшейся части этого дерева каждая вершина имеет степень
1 плюс число её вхождений в оставшемся коде: каждое такое вхождение
прикрепляет к ней один удаляемый ранее лист, а одно ребро связывает её
с более поздней вершиной либо является последним ребром.
Следовательно, на каждом шаге выбранная \(x\) действительно есть
наименьший лист оставшегося дерева. Повторное кодирование даёт \(a_t\)
на каждом шаге, то есть исходное слово.
\end{proof}
\begin{theorem}[Биекция Прюфера и формула Кэли]
Кодирование задаёт биекцию между деревьями на \([n]\) и словами
длины \(n-2\) над алфавитом \([n]\). Поэтому число помеченных деревьев
при \(n\ge2\) равно \(n^{n-2}\).
\end{theorem}
\begin{proof}
Декодирование любого слова даёт дерево с этим кодом.
Если начать с дерева, лемма о частотах показывает, что декодер
на каждом шаге выбирает именно лист, удалённый кодером, и возвращает
все исходные рёбра. Таким образом, алгоритмы взаимно обратны.
У каждой из \(n-2\) позиций \(n\) независимых вариантов, всего \(n^{n-2}\).
\end{proof}
При \(n=2\) код пуст и дерево одно; при \(n=1\) дерево также одно,
этот случай рассматривается отдельно.
\begin{example}
Для звезды с центром 4 на \(\{1,2,3,4,5\}\) код равен \((4,4,4)\).
Декодер последовательно присоединяет к 4 вершины \(1,2,3\), затем
соединяет последние \(4,5\). Частота 4 равна \(3=d(4)-1\).
\end{example}

\ticket{Эйлеровы графы}
\label{ticket:9}
\source{Лекции, с. 28–29; строгие условия связности и детали индукции восстановлены составителем; Карпов, §1.9}
\begin{definition}[Эйлеров путь и цикл]
Эйлеров путь проходит по каждому ребру графа ровно один раз.
Если он замкнут, это эйлеров цикл; граф с таким циклом называется эйлеровым.
Вершины могут повторяться. В орграфе дуги проходят только по направлению.
Далее существенная часть — вершины ненулевой степени со всеми рёбрами.
В лекциях исходный граф предполагается связным; отделение существенной
части позволяет также учитывать изолированные вершины.
\end{definition}
\begin{theorem}[Критерии Эйлера]
Для графа с хотя бы одним ребром необходима связность существенной части
после забывания направлений. При этом условии верны следующие критерии:
\begin{center}
\begin{tabular}{p{.13\linewidth}|p{.34\linewidth}|p{.39\linewidth}}
 & Замкнутый обход & Открытый обход с концами \(s\ne t\)\\\hline
Граф & Все степени чётны & Ровно две нечётные вершины \(s,t\)\\[3pt]
Орграф & \(d^+(v)=d^-(v)\) для всех \(v\) &
\(d^-(s)-d^+(s)=1\), \(d^+(t)-d^-(t)=1\); у остальных баланс нулевой
\end{tabular}
\end{center}
\end{theorem}
\begin{proof}
Необходимость: у внутренних посещений вход и выход образуют пару.
У начала открытого пути остаётся один лишний выход, у конца — вход.
Один обход всех рёбер также даёт требуемую связность.

\textit{Неориентированный цикл: лекционная индукция по числу рёбер.}
В графе с положительными чётными степенями есть простой цикл:
продолжаем простой путь до первой вынужденной встречи с прежней вершиной.
Удаляем рёбра этого цикла \(C\). В каждой оставшейся существенной
компоненте все степени чётны и рёбер меньше; по индукции есть эйлеров
цикл. Каждая такая компонента имеет общую вершину с \(C\): иначе исходная
связность нарушалась бы (все удалённые рёбра имеют концы на \(C\)).
Начинаем обход \(C\), а при посещении общей вершины вставляем обход
соответствующей компоненты и возвращаемся к обходу \(C\).
Получено прохождение всех рёбер ровно один раз. Если остаточных рёбер нет,
сам \(C\) завершает доказательство; это также база индукции.

\textit{Ориентированный цикл.} Пока есть дуги, в сбалансированном орграфе
есть направленный цикл. Действительно, после входа в вершину у неё есть
исходящая дуга, а бесконечный выбор продолжений в конечном множестве
даёт цикл. Удаление дуг простого направленного цикла сохраняет баланс.
Остаточные компоненты графа с забытой ориентацией сбалансированы;
по индукции у них есть направленные эйлеровы циклы. Склейка с исходным
циклом проводится так же. Здесь сильная связность не предполагается:
она сама следует из полученного обхода всех ненулевых вершин.

\textit{Открытый обход.} В неориентированном случае добавим ребро \(st\),
в ориентированном — дугу \(t\to s\). Получится граф с условием для цикла.
Строим этот цикл, затем разрезаем его запись по добавленному ребру/дуге.
Остаётся эйлеров путь из \(s\) в \(t\). Временные кратные рёбра допустимы:
доказательство индукцией различает рёбра, а не только их концы.
\end{proof}
\begin{cor}[Разбиение на пути]
Если связный неориентированный граф имеет \(2k>0\) нечётных вершин,
его рёбра разбиваются на \(k\) непересекающихся по рёбрам открытых путей.
Меньшим числом путей обойти все рёбра невозможно.
\end{cor}
\begin{proof}
Попарно соединим нечётные вершины \(k\) временными рёбрами.
Полученный связный мультиграф имеет все степени чётными. Разрезая его
эйлеров цикл по временным рёбрам, получаем \(k\) путей исходного графа.
Временные рёбра не могут следовать подряд, поскольку имеют разные концы;
поэтому все полученные пути непустые. Каждая нечётная вершина должна
быть концом нечётного числа путей разбиения. У \(r\) путей только \(2r\)
концов, следовательно, \(r\ge k\).
\end{proof}
Если все степени чётны и есть ребро, достаточно одного замкнутого пути.
Для одноточечного графа без рёбер нулевой путь считается эйлеровым.
Петли, если они специально разрешены, дают две инцидентности в неориентированном
случае и по одному входу и выходу в ориентированном; критерии сохраняются.

\ticket{Гамильтоновы графы. Теорема Хватала}
\label{ticket:10}
\source{Лекции, с. 30–31; заключительный подсчёт восстановлен составителем; Карпов, §3.1}
\begin{definition}[Гамильтонов граф]
Гамильтонов путь — простой путь, содержащий все вершины графа.
Гамильтонов цикл — простой цикл, содержащий все вершины.
Граф называется гамильтоновым, если имеет гамильтонов цикл; далее \(n\ge3\).
Степенная последовательность — степени в неубывающем порядке
\(d_1\le\cdots\le d_n\).
\end{definition}
\begin{theorem}[Хватал]
Пусть \(n\ge3\), \(d_1\le\cdots\le d_n\). Если для каждого целого
\(1\le i<n/2\) выполнено
\[
\boxed{d_i\le i\quad\Longrightarrow\quad d_{n-i}\ge n-i,}
\]
то граф имеет гамильтонов цикл.
\end{theorem}
\begin{proof}
Предположим противное. Добавление рёбер не уменьшает ни одну
из упорядоченных степеней, поэтому сохраняет условие: если новая
\(d_i\le i\), то и старая \(d_i\le i\), а новая \(d_{n-i}\) не меньше старой.
Добавим к исходному графу как можно больше рёбер, сохраняя негамильтоновость;
получим \(H\). Он не полный, а добавление любого отсутствующего ребра
делает его гамильтоновым. В гамильтоновом цикле \(H+uv\) обязательно
используется новое ребро, иначе \(H\) уже гамильтонов.
Удалив его, получаем гамильтонов путь
\(x_1=u,x_2,\ldots,x_n=v\) в \(H\).

Среди несмежных \(u,v\) выберем пару с максимальной суммой степеней,
\(h=d_H(u)\le d_H(v)=q\). Рассмотрим
\[
S=\{i:x_i v\in E(H)\},\qquad
T=\{i:u x_{i+1}\in E(H)\}\subseteq\{1,\ldots,n-1\}.
\]
Имеем \(|S|=q\), \(|T|=h\). Если \(i\in S\cap T\), возникает цикл
\[
u,x_2,\ldots,x_i,v,x_{n-1},\ldots,x_{i+1},u,
\]
проходящий через все вершины. Крайние вырожденные индексы исключены
несмежностью \(u,v\). Следовательно, \(S\cap T=\varnothing\), и
\(h+q\le n-1\), в частности \(h<n/2\).

Каждый несосед \(v\), отличный от \(v\), имеет степень не больше \(h\),
иначе выбранная сумма не максимальна. Таких вершин \(n-1-q\ge h\).
Каждый несосед \(u\) имеет степень не больше \(q\le n-h-1\).
Вместе с \(u\) это \(n-h\) вершин степени не больше \(n-h-1\)
(\(h\le n-h-1\)). Если \(h\ge1\), то
\[
d_h(H)\le h,\qquad d_{n-h}(H)\le n-h-1,
\]
что противоречит условию Хватала. Если \(h=0\), вершина \(u\) изолирована,
а все остальные степени не больше \(n-2\). Тогда
\(d_1(H)=0\le1\), \(d_{n-1}(H)\le n-2\), снова противоречие при \(i=1<n/2\).
\end{proof}
\begin{center}
\begin{tikzpicture}[>=Stealth,scale=.9,every node/.style={circle,fill=cbBlue,inner sep=2pt}]
\foreach \x/\name in {0/u,1.1/a,2.2/b,3.3/c,4.4/d,5.5/v}{\node (\name) at (\x,0) {};}
\draw (u)--(a)--(b)--(c)--(d)--(v);
\draw[cbOrange,thick] (u) to[bend left=35] (c);
\draw[cbGreen,thick] (b) to[bend right=35] (v);
\node[fill=none,below] at (u) {$u=x_1$};\node[fill=none,below] at (b) {$x_i$};
\node[fill=none,below] at (c) {$x_{i+1}$};\node[fill=none,below] at (v) {$v=x_n$};
\end{tikzpicture}
\end{center}
Два выделенных ребра при \(i\in S\cap T\) замыкают путь в гамильтонов цикл.

\textbf{Книжное дополнение для следующего билета.}
\begin{lemma}[Добавление ребра]\label{lem:closure}
Если \(u,v\) несмежны и \(d(u)+d(v)\ge n\), то \(G\) гамильтонов
тогда и только тогда, когда \(G+uv\) гамильтонов.
\end{lemma}
\begin{proof}
Только обратное направление требует проверки. Если цикл использует
\(uv\), удаляем его и получаем гамильтонов путь. Для множеств \(S,T\)
выше сумма мощностей не меньше \(n\), поэтому они пересекаются в
множестве из \(n-1\) индексов. Та же перестройка даёт цикл в \(G\).
\end{proof}
\begin{definition}[Замыкание; дополнение по книге]
Последовательно добавляя рёбра между несмежными вершинами с суммой
степеней не меньше \(n\), получаем замыкание \(\cl(G)\).
Оно одновременно гамильтоново с \(G\) и не зависит от порядка добавлений.
\end{definition}
\begin{proof}[Обоснование]
Первое следует из леммы. Для двух завершённых процессов \(H_1,H_2\)
индукцией по добавлениям первого покажем \(H_1\subseteq H_2\).
Если предыдущие рёбра уже есть в \(H_2\), степени концов следующего
там не меньше, чем при его добавлении. Отсутствие ребра в \(H_2\)
противоречило бы завершённости второго процесса. Симметрично получаем
обратное включение. В частности, \(\cl(G)=K_n\) влечёт гамильтоновость.
\end{proof}
\textbf{Дополнительное свойство степенного критерия (по книге).}
\begin{cor}[Точность критерия как условия на степени]
Если неубывающая последовательность \(0\le a_1\le\cdots\le a_n\le n-1\)
нарушает условие Хватала при \(i=h\), существует негамильтонов граф,
степени которого покоординатно не меньше \(a_1,\ldots,a_n\).
\end{cor}
\begin{proof}
Возьмём три множества \(A,B,D\) размеров \(h,h,n-2h\).
В \(B\cup D\) проведём все рёбра; соединим каждую вершину \(A\)
со всеми вершинами \(B\), других рёбер у \(A\) нет.
Степени равны \(h\) на \(A\), \(n-h-1\) на \(D\), \(n-1\) на \(B\),
поэтому они мажорируют заданную последовательность.
После удаления \(B\) остаются \(h+1\) компонент: отдельные вершины \(A\)
и непустая клика \(D\). Удаление \(h\) вершин из гамильтонова цикла
оставляет не более \(h\) непустых отрезков, значит, в гамильтоновом графе
компонент после такого удаления не больше \(h\). Противоречие.
\end{proof}
Критерий достаточен для конкретного графа, но не необходим: например,
цикл \(C_5\) гамильтонов, а при \(i=2\) его степени нарушают условие.

\ticket{Теоремы Оре и Дирака}
\label{ticket:11}
\begin{definition}
В этом билете граф простой, \(n\ge3\),
\(d_1\le\cdots\le d_n\) — упорядоченные степени.
Гамильтоновость означает наличие простого цикла через все вершины.
\end{definition}
\begin{theorem}[Оре]
Если для любых двух несмежных вершин \(u,v\)
\(d(u)+d(v)\ge n\), граф гамильтонов.
\end{theorem}
\begin{proof}
Выведем утверждение из критерия Хватала.


Предположим, что критерий Хватала нарушен при \(1\le i<n/2\):
\(d_i\le i\), \(d_{n-i}\le n-i-1\). Пусть \(T\) — первые \(n-i\)
вершин упорядоченного списка, \(S\subset T\) — первые \(i\).
Для \(x\in S\) и \(y\in T\setminus\{x\}\) сумма степеней не превосходит
\(i+(n-i-1)=n-1\). По условию Оре \(x,y\) обязательно смежны.
Отсюда \(i\ge d(x)\ge n-i-1\ge i\), где последнее неравенство
следует из \(n>2i\). Все неравенства должны быть равенствами:
\(n=2i+1\), \(d(x)=i\), и у \(x\) нет соседей вне \(T\).
В \(T\setminus S\) осталась ровно одна вершина \(z\).
Она смежна со всеми \(i\) вершинами \(S\) и имеет степень не больше \(i\),
так что тоже не имеет соседей вне \(T\).
Получили компоненту \(T=K_{i+1}\); дополнение размера \(i\) непусто
и его вершины имеют степень не больше \(i-1\).
Выбрав \(x\in T\), \(w\notin T\), получаем несмежную пару с суммой
степеней не больше \(2i-1<n\), противоречие Оре.
Следовательно, условие Хватала выполнено и граф гамильтонов.
\end{proof}
Есть и более короткое доказательство через замыкание из билета 10:
любую недостающую пару можно дополнить ребром, условие Оре при этом
сохраняется, в итоге получаем \(K_n\). Лемма о добавлении ребра
сохраняет гамильтоновость в обе стороны.
\begin{theorem}[Дирак]
Если \(\delta(G)\ge n/2\), то граф гамильтонов.
\end{theorem}
\begin{proof}
Для любых несмежных \(u,v\) имеем \(d(u)+d(v)\ge2\delta(G)\ge n\).
Применима теорема Оре. Непосредственно условие Хватала тоже выполнено:
при \(i<n/2\) имеем \(d_i\ge\delta(G)>i\), поэтому посылка импликации ложна.
\end{proof}
\begin{cor}[Версии для пути]
Если \(n\ge2\) и сумма степеней любой несмежной пары не меньше \(n-1\),
есть гамильтонов путь. Достаточно также условия \(\delta(G)\ge(n-1)/2\).
\end{cor}
\begin{proof}
Добавим универсальную вершину \(w\). В новом графе \(n+1\) вершина,
а степени старой несмежной пары в сумме не меньше \((n-1)+2=n+1\).
У \(w\) несмежных пар нет, поэтому применима теорема Оре.
Удаление \(w\) из гамильтонова цикла оставляет гамильтонов путь.
Условие на минимальную степень влечёт условие на сумму.
\end{proof}
\begin{example}[Оценку Дирака нельзя понизить на единицу]
При \(n=2r\), \(r\ge2\), граф \(K_{r-1,r+1}\) имеет минимальную степень
\(r-1\), но не имеет гамильтонова цикла: в двудольном цикле числа
вершин обеих долей равны. Условия теорем являются достаточными;
их нарушение само по себе не доказывает негамильтоновость.
\end{example}

\ticket{Гамильтоновы турниры. Теоремы Редеи и Камиона}
\label{ticket:12}
\begin{definition}[Турнир]
Турнир — орграф, в котором между любыми двумя различными вершинами
есть ровно одна дуга одного из двух направлений; петель нет.
Гамильтоновы путь и цикл проходят через все вершины по направлениям дуг.
\end{definition}
\begin{lemma}[Вставка в путь]
Если \(v_1\to\cdots\to v_k\) — путь в турнире и \(x\) — новая вершина,
то её можно вставить в этот путь с сохранением порядка старых вершин.
\end{lemma}
\begin{proof}
Если \(x\to v_1\), добавим \(x\) в начало; если \(v_k\to x\), в конец.
Иначе \(v_1\to x\) и \(x\to v_k\). Пусть \(j\) — первый индекс,
для которого \(x\to v_j\). Тогда \(j\ge2\), а \(v_{j-1}\to x\),
поэтому вставка между \(v_{j-1}\) и \(v_j\) корректна.
\end{proof}
\begin{theorem}[Редеи]
В любом непустом турнире существует гамильтонов путь.
\end{theorem}
\begin{proof}
Начинаем с одной вершины и вставляем остальные по лемме.
После каждой вставки путь содержит все уже обработанные вершины;
через конечное число шагов он содержит весь турнир.
\end{proof}
\begin{lemma}[Вставка в цикл]
Если вершина \(x\) вне направленного цикла имеет как входящие,
так и исходящие дуги к его вершинам, её можно вставить в цикл.
\end{lemma}
\begin{proof}
При циклическом просмотре вершин есть переход от типа
\(v_i\to x\) к типу \(x\to v_{i+1}\), поскольку оба типа присутствуют.
Заменим дугу \(v_i\to v_{i+1}\) двумя дугами через \(x\).
\end{proof}
\begin{graphillustration}{Вставка $x$ заменяет одну дугу двумя. На схеме цикла показан пример с тремя исходными вершинами; на пути возможна также вставка на конце.}{fig:12-insert}
\begin{tikzpicture}[font=\small]
\begin{scope}[shift={(0,0)}]
\node[gi vertex] (p) at (0,0) {$v_{j-1}$};
\node[gi vertex] (q) at (3.5,0) {$v_j$};
\node[gi selected] (x) at (1.75,1.3) {$x$};
\draw[gi removed,-{Stealth}] (p) -- (q);
\draw[gi arrow] (p) -- (x);
\draw[gi arrow] (x) -- (q);
\node[gi note] at (-0.7,0) {$\cdots$};
\node[gi note] at (4.2,0) {$\cdots$};
\node[gi note] at (1.75,-1.9) {Вставка внутри пути};
\end{scope}
\begin{scope}[shift={(6.5,0)}]
\node[gi vertex] (p) at (0,0) {$v_i$};
\node[gi vertex] (q) at (3.5,0) {$v_{i+1}$};
\node[gi selected] (x) at (1.75,1.3) {$x$};
\draw[gi removed,-{Stealth}] (p) -- (q);
\draw[gi arrow] (p) -- (x);
\draw[gi arrow] (x) -- (q);
\node[gi vertex] (w) at (1.75,-1.1) {$w$};
\draw[gi arrow] (q) -- (w);
\draw[gi arrow] (w) -- (p);
\node[gi note] at (1.75,-1.9) {Вставка в цикл};
\end{scope}
\end{tikzpicture}
\end{graphillustration}


\begin{theorem}[Камион; критерий Редеи–Камиона]
Турнир с \(n\ge3\) имеет гамильтонов цикл тогда и только тогда,
когда он сильно связен.
\end{theorem}
\begin{proof}
Гамильтонов цикл даёт направленный путь между любой упорядоченной парой,
значит, сильная связность необходима.
Пусть турнир сильно связен. В нём есть направленный цикл:
возьмём дугу \(u\to v\) и простой обратный путь из \(v\) в \(u\);
его длина не меньше двух, поскольку противоположной дуги нет.
Выберем простой цикл \(C\) максимальной длины.
Если \(C\) не гамильтонов, по лемме каждая внешняя вершина либо
получает дуги от всех вершин \(C\), либо посылает дуги ко всем ним.
Обозначим эти множества \(A\) и \(B\) соответственно.

Оба множества непусты: при \(A=\varnothing\) из цикла нельзя выйти,
а при \(B=\varnothing\) нельзя в него попасть, вопреки сильной связности.
Есть дуга \(a\to b\), \(a\in A,b\in B\): иначе все дуги между ними
направлены из \(B\) в \(A\) и из \(A\) нельзя достичь цикла.
Теперь любую дугу \(v_i\to v_{i+1}\) цикла можно заменить путём
\(v_i\to a\to b\to v_{i+1}\), увеличив цикл на две вершины.
Противоречие максимальности. Поэтому \(C\) гамильтонов.
\end{proof}

\begin{graphillustration}{Шаг доказательства Камиона: внешним путём $v_i\to a\to b\to v_{i+1}$ заменяем дугу цикла. Схема показывает только используемые дуги турнира.}{fig:12-camion}
\begin{tikzpicture}[font=\small]
\node[gi vertex] (vi) at (0,0) {$v_i$};
\node[gi vertex] (vj) at (4.5,0) {$v_{i+1}$};
\node[gi vertex] (w) at (2.25,-1.3) {$w$};
\node[gi vertex] (a) at (1,1.5) {$a\in A$};
\node[gi vertex] (b) at (3.5,1.5) {$b\in B$};
\draw[gi arrow] (vi) -- (a);
\draw[gi arrow] (a) -- (b);
\draw[gi arrow] (b) -- (vj);
\draw[gi arrow] (vj) -- (w);
\draw[gi arrow] (w) -- (vi);
\draw[gi removed,-{Stealth}] (vi) -- (vj);
\node[gi note] at (7.7,0.5) {$C\to A,\quad A\ni a\to b\in B,\quad B\to C$\\[5pt]Цикл удлиняется на две вершины.};
\end{tikzpicture}
\end{graphillustration}


\ticket{Плоские графы, формула Эйлера и гомеоморфизм}
\label{ticket:13}
% Чуть компактнее вертикальные интервалы, чтобы завершить билет на одной странице.
\begingroup
\setlength{\parskip}{2pt}
\tcbset{before skip=4pt,after skip=4pt,top=5pt,bottom=5pt}
\source{Лекции, с. 32–37; Карпов, §6.1, §6.3; детали индукции восстановлены составителем}
\begin{definition}[Укладка на поверхности и грани]
Пусть \(M\) — поверхность. Укладка графа \(G\) на \(M\) состоит
из инъекции \(\varphi_V:V\to M\) и простых дуг
\(\varphi_e:[0,1]\to M\), \(e=uv\), с концами
\(\varphi_e(0)=\varphi_V(u)\), \(\varphi_e(1)=\varphi_V(v)\).
Внутренность дуги не содержит образов вершин, а две дуги могут встречаться
только в образах общих концов соответствующих рёбер.
Планарный граф допускает укладку на плоскости; плоский граф — граф
с уже выбранной укладкой. Грани — компоненты дополнения к рисунку;
на плоскости ровно одна грань неограниченна и называется внешней.
Число граней обозначается \(f\). Можно работать с кусочно-линейными дугами.
\end{definition}
\begin{lemma}[Удаление рёбер на рисунке]
Ребро цикла граничит с двумя различными гранями; его удаление объединяет
эти грани и уменьшает \(f\) на 1. Висячее ребро не разделяет грань;
его удаление вместе с листом сохраняет число граней.
\end{lemma}
\begin{proof}
Простой цикл через ребро — кривая Жордана. Две стороны внутренней
точки ребра лежат по разные стороны этой кривой и потому не могут
принадлежать одной грани. Малый открытый коридор на месте удалённого
ребра соединяет именно эти две области; других изменений дополнения нет.
У висячего ребра обе стороны соединяются обходом свободного конца
в малой окрестности, не содержащей других рёбер. Оно лишь образует
разрез внутри одной области, а не отделяет новую область.
\end{proof}
\begin{theorem}[Формула Эйлера]
Для плоского графа с \(c\) компонентами
\[
 \boxed{n-m+f=1+c.}
\]
В частности, для связного непустого графа \(n-m+f=2\).
\end{theorem}
\begin{proof}
Сначала пусть граф связен. Строим его рисунок из одной вершины.
По билету 6 есть остовное дерево; укореним его и добавим остальные
вершины по одной вместе с ребром к уже построенному родителю.
Числа \(n,m\) возрастают на 1, число граней не меняется: висячее ребро
грань не разделяет. Затем добавляем остальные рёбра в их исходных дугах.
Концы очередного ребра уже соединены остовом, поэтому новое ребро
замыкает простой цикл и делит одну грань на две по теореме Жордана.
Числа \(m,f\) возрастают на 1. Оба шага сохраняют \(n-m+f\).
В базе \(n=1,m=0,f=1\), значит, для связного графа \(n-m+f=2\).

\textit{Обобщение составителя.} Для несвязного графа можно соединить разные компоненты рёбрами без
пересечений: берём дугу между рисунками двух компонент, и участок между
последовательными касаниями разных компонент, внутренность которого
не встречает граф. Его концы при необходимости доводим вдоль границ
грани до вершин. Новое ребро — мост и не создаёт грани: его стороны
соединяются обходом одной из прежде раздельных компонент.
После \(c-1\) добавлений получим связный граф с теми же \(n,f\)
и \(m+c-1\) рёбрами. Отсюда \(n-(m+c-1)+f=2\).
Для пустого графа \(n=m=c=0,f=1\), равенство также верно.
\end{proof}
\begin{definition}[Подразбиение и гомеоморфные графы]
Элементарные операции: заменить ребро \(uv\) на \(ut,tv\) с новой
вершиной \(t\); обратно, удалить вершину степени 2, склеив две её дуги
в одно ребро (при необходимости разрешая промежуточные кратности).
Графы гомеоморфны, если связаны последовательностью таких операций
и изоморфизмов. Эквивалентно, они имеют изоморфные подразбиения.
Подразбиение — замена рёбер простыми путями с новыми внутренними
вершинами; у разных рёбер эти вершины различны. Это комбинаторное понятие;
гомеоморфизм плоскости — непрерывная биекция с непрерывной обратной.
\end{definition}
\begin{theorem}[Инвариантность планарности]
Граф планарен тогда и только тогда, когда планарно любое его подразбиение.
Следовательно, гомеоморфные графы одновременно планарны.
\end{theorem}
\begin{proof}
При подразделении отмечаем новые вершины на дугах рёбер. Обратно,
склеиваем последовательные дуги каждого заменяющего пути в одну простую
дугу; других пересечений при этом не возникает. Гомеоморфные графы
сравниваются через общее подразбиение.
\end{proof}

\endgroup

\ticket{Укладка на сфере и выбор внешней грани}
\label{ticket:14}
\source{Лекции, с. 34–35, 38–39; Карпов, §6.1.3, §6.2.1; координатные формулы — пояснение составителя}
\begin{definition}[Сферическая укладка]
Укладка на сфере определяется так же, как на плоскости: вершины — точки
сферы, рёбра — простые дуги без посторонних пересечений.
Все грани на сфере равноправны, выделенной внешней грани нет.
\end{definition}
\begin{lemma}[Стереографическая проекция]
Сфера с удалённой точкой гомеоморфна плоскости.
\end{lemma}
\begin{proof}
Для единичной сферы \(x^2+y^2+z^2=1\) и северного полюса \(N=(0,0,1)\)
проекция на плоскость \(z=0\) задаётся формулой
\[
 (x,y,z)\longmapsto\left(\frac{x}{1-z},\frac{y}{1-z}\right).
\]
Её обратное отображение для \(s=u^2+v^2\):
\[
 (u,v)\longmapsto\left(\frac{2u}{s+1},\frac{2v}{s+1},\frac{s-1}{s+1}\right).
\]
Подстановка проверяет, что образы лежат на сфере и отображения взаимно обратны.
Оба непрерывны, поскольку знаменатели не обращаются в нуль на своих областях.
Для любой другой удаляемой точки предварительно повернём сферу.
\end{proof}
\begin{theorem}[Плоскость и сфера]
Граф допускает укладку на плоскости тогда и только тогда,
когда допускает укладку на сфере.
\end{theorem}
\begin{proof}
Обратная стереографическая проекция переносит плоскую укладку на сферу.
Для обратного перехода выбираем на сфере точку, не лежащую на рисунке,
и проектируем из неё. Такая точка есть: конечное объединение
кусочно-гладких дуг и точек не покрывает сферу.
Гомеоморфизм сохраняет инъективность каждой дуги и отсутствие пересечений.
\end{proof}
\begin{theorem}[Выбор внешней грани]
Любую грань данной сферической укладки можно сделать внешней в некоторой
плоской укладке того же графа.
\end{theorem}
\begin{proof}
Выберем полюс проекции внутри нужной грани. При приближении к полюсу
образ уходит на бесконечность, поэтому образ этой грани неограничен.
Остальные грани ограничены: замыкание каждой не содержит выбранного
полюса и является компактным подмножеством области проекции.
Таким образом, выбранная грань становится внешней.
\end{proof}
\begin{cor}[Вершина и ребро на внешней грани]
Для любой выбранной вершины или любого выбранного ребра планарного
графа существует укладка, в которой этот объект лежит на границе внешней грани.
\end{cor}
\begin{proof}
Выбираем инцидентную объекту грань и делаем её внешней по теореме.
Для вершины такая грань есть в любом свободном секторе её окрестности;
для изолированной вершины — в окружающем её дополнении.
Грань, инцидентная внутренней точке ребра, инцидентна всей его дуге,
поскольку внутри ребра нет вершин. Для моста обе стороны принадлежат
одной грани, что не меняет рассуждения.
\end{proof}

\ticket{Непланарность полного графа и полного двудольного графа}
\label{ticket:15}
\begin{definition}[Длина границы грани]
Длина границы — число сторон рёбер при обходе грани с учётом кратности.
Мост учитывается дважды на границе одной грани. Поэтому
\(\sum_F \ell(F)=2m\): у каждого ребра две стороны.
\end{definition}
\begin{graphillustration}{Пример обхода границ: мост $cv$ встречается на границе внешней грани дважды — по одной стороне в каждом направлении.}{fig:15-boundaries}
\begin{tikzpicture}[font=\small]
\node[gi vertex] (a) at (0,0) {$a$};
\node[gi vertex] (b) at (2,0) {$b$};
\node[gi vertex] (c) at (1,1.7) {$c$};
\node[gi vertex] (v) at (1,3) {$v$};
\draw[gi edge] (a) -- (b);
\draw[gi edge] (b) -- (c);
\draw[gi edge] (c) -- (a);
\draw[gi edge] (c) -- (v);
\draw[gi flow] (.8,2)--(.8,2.65);\draw[gi flow] (1.2,2.65)--(1.2,2);
\node[gi note] at (1,0.65) {$F_1$};
\node[gi note] at (3,1.8) {$F_0$};
\node[gi note] at (6.3,1.2) {$\ell(F_1)=3$\\$\ell(F_0)=3+2=5$\\[5pt]$\ell(F_0)+\ell(F_1)=8=2m$};
\end{tikzpicture}
\end{graphillustration}


\begin{lemma}[Оценки числа рёбер]
Для простого планарного графа с \(n\ge3\)
\[
 m\le3n-6.
\]
Если граф двудолен, то \(m\le2n-4\).
\end{lemma}
\begin{proof}
Пусть сначала граф связен. Если это дерево, \(m=n-1\), и обе оценки
следуют из \(n\ge3\). Иначе любая граница грани имеет длину не менее 3:
длина 1 требовала бы петли, длина 2 — кратных рёбер или единственного
ребра, дважды обходимого как весь связный граф. Последнее невозможно
при \(n\ge3\). В двудольном графе все замкнутые маршруты имеют чётную
длину, так что длина грани не менее 4.

В первом случае \(3f\le2m\), и по Эйлеру
\(2=n-m+f\le n-m/3\), откуда \(m\le3n-6\).
Во втором \(4f\le2m\), следовательно,
\(2\le n-m/2\), что даёт \(m\le2n-4\).
Если граф несвязен, соединим его компоненты непересекающимися мостами.
Простота сохраняется; в двудольном случае при необходимости меняем
названия долей в одной из соединяемых компонент. Получим связный
простой (соответственно двудольный) планарный граф на тех же вершинах
с не меньшим числом рёбер. Его оценка влечёт исходную.
\end{proof}
\begin{theorem}
Графы \(K_5\) и \(K_{3,3}\) непланарны.
\end{theorem}
\begin{proof}
У \(K_5\) имеем \(n=5,m=10\), но для планарного графа
\(m\le3\cdot5-6=9\), противоречие.
У двудольного \(K_{3,3}\) имеем \(n=6,m=9\), но
\(m\le2\cdot6-4=8\), снова противоречие.
\end{proof}
\begin{center}
\begin{tikzpicture}[every node/.style={circle,fill=cbBlue,inner sep=2pt}]
\foreach \i in {1,...,5}{\node (p\i) at ({90+72*(\i-1)}:1.15) {};}
\foreach \i in {1,...,5}{\foreach \j in {1,...,5}{\ifnum\i<\j\draw (p\i)--(p\j);\fi}}
\node[fill=none,rectangle] at (0,-1.6) {$K_5$};
\foreach \i in {1,2,3}{\node (a\i) at (3,{2-\i}) {};\node (b\i) at (5,{2-\i}) {};}
\foreach \i in {1,2,3}{\foreach \j in {1,2,3}{\draw (a\i)--(b\j);}}
\node[fill=none,rectangle] at (4,-1.6) {$K_{3,3}$};
\end{tikzpicture}
\end{center}
Пересечения на этих рисунках сами по себе не доказывают непланарность:
доказательство должно исключать \emph{любую} укладку, что и делают оценки.
\begin{cor}
Ни один планарный граф не содержит подграфа, являющегося подразбиением
\(K_5\) или \(K_{3,3}\).
\end{cor}
\begin{proof}
Подграф планарного графа планарен: достаточно удалить лишние дуги и точки
из укладки. Планарность подразбиения равносильна планарности исходного
графа (билет 13), тогда \(K_5\) или \(K_{3,3}\) был бы планарен.
\end{proof}

\ticket{Планарность и компоненты двусвязности}
\label{ticket:16}
% Чуть компактнее вертикальные интервалы, чтобы завершить билет на одной странице.
\begingroup
\setlength{\parskip}{2pt}
\tcbset{before skip=4pt,after skip=4pt,top=5pt,bottom=5pt}
\source{Лекции, с. 39–41; Карпов, §6.4; детали склейки раскрыты составителем}
\begin{definition}
Рёберные компоненты двусвязности — компоненты после удаления всех мостов
(билет 4). Вершинные компоненты двусвязности — блоки (билет 5), включая
мосты с концами и изолированные вершины.
\end{definition}
\begin{lemma}[Склейка в вершине]
Если два планарных графа имеют ровно одну общую вершину и не имеют
общих рёбер, то их объединение планарно.
\end{lemma}
\begin{proof}
Рисунок конечного графа — конечное объединение компактных дуг и точек,
поэтому ограничен. Переносом и гомотетией его можно поместить в любой
заданный открытый диск, сохраняя отсутствие пересечений.
Выберем в каждом рисунке общую вершину \(a\) на внешней грани
(билет 14). Укладку каждого графа можно поместить в замкнутый
топологический диск, граница которого проходит через \(a\) и больше
нигде не касается рисунка. Для этого в секторе внешней грани у \(a\)
проведём два близких непересекающихся выхода к большой окружающей
окружности и соединим их по окружности; полученная простая замкнутая
кривая ограничивает диск с рисунком внутри. Для изолированной вершины
достаточно любого касающегося её диска.
Деформируем диски в два диска, касающиеся только в \(a\), и отождествим
вершины. Внутренности дисков раздельны, новых пересечений нет.
Это также описывает локальную склейку конечного числа рисунков в разных секторах.
\end{proof}
\begin{lemma}[Добавление моста между рисунками]
Два непересекающихся планарных графа можно соединить ребром
между любыми выбранными вершинами с сохранением планарности.
\end{lemma}
\begin{proof}
Делаем обе выбранные вершины внешними, помещаем рисунки в раздельные
диски с этими вершинами на границе и проводим соединяющую дугу
снаружи дисков. Её внутренность не пересекает рисунков.
\end{proof}
\begin{theorem}[Планарность по рёберным компонентам]
Граф планарен тогда и только тогда, когда планарны все его
рёберные компоненты двусвязности.
\end{theorem}
\begin{proof}
Необходимость — наследование планарности подграфами.
После сжатия компонент остаётся лес мостов: цикл в нём дал бы обход
одного из мостов, что противоречит определению моста.
Для каждого дерева этого леса выбираем корень и последовательно
присоединяем рисунки дочерних компонент мостами по предыдущей лемме.
Каждый шаг соединяет два до того раздельных рисунка, поэтому пересечений нет.
Разные деревья леса размещаем отдельно.
\end{proof}
\begin{theorem}[Планарность по блокам]
Граф планарен тогда и только тогда, когда каждый его блок планарен.
\end{theorem}
\begin{proof}
Необходимость следует из планарности подграфов.
Для достаточности сначала рассматриваем связный граф и индукцию
по числу блоков. Если блок один, утверждение дано.
Иначе по дереву блоков (билет 5) существует крайний блок \(B\),
соединённый с остальной частью графа только через одну точку сочленения \(a\).
Остальная часть имеет меньше блоков и планарна по индукции.
Склеиваем её рисунок с рисунком \(B\) в вершине \(a\) по лемме.
Для несвязного графа рисуем его компоненты в непересекающихся дисках.
Изолированные вершины и блоки-мосты заведомо планарны.
\end{proof}
\begin{center}
\begin{tikzpicture}[scale=0.8,every node/.style={circle,fill=cbBlue,inner sep=2pt}]
\node (a) at (0,0) {};\node (b) at (-1,1) {};\node (c) at (-1,-1) {};
\node (d) at (1,1) {};\node (e) at (1,-1) {};
\draw (a)--(b)--(c)--(a)--(d)--(e)--(a);
\node[rectangle,fill=none,above] at (a) {$a$};
\node (f) at (2.7,-1) {};\draw (e)--(f);
\end{tikzpicture}
\end{center}
Здесь два треугольника — разные вершинные блоки, но лежат в одной
рёберной компоненте двусвязности. Последнее ребро — мост.

\endgroup

\ticket{Теорема Понтрягина–Куратовского}
\label{ticket:17}
\begingroup
\setlength{\parskip}{2pt}
\tcbset{before skip=4pt,after skip=4pt,top=4pt,bottom=4pt}
\source{Лекции, с. 41–45; схема через цикл и его части; минимальность и локальные конфигурации раскрыты составителем}
\begin{definition}[Топологический подграф]
Подграф, гомеоморфный \(H\), — подразбиение \(H\): его рёбра заменены
простыми путями с различными новыми внутренними вершинами.
Вершины, соответствующие вершинам \(H\), называются главными.
\end{definition}
\begin{theorem}[Понтрягин–Куратовский]
Граф планарен тогда и только тогда, когда не содержит подграфа,
гомеоморфного \(K_5\) или \(K_{3,3}\).
\end{theorem}
\begin{proof}[Необходимость]
Подграф планарного графа планарен, подразделение и обратное склеивание
рёбер сохраняют планарность, а \(K_5\), \(K_{3,3}\) непланарны по билету 15.
\end{proof}
\subsection*{Лекционная схема достаточности}
В вопросе требуется \emph{схема}. Ниже сохранён ход лекций:
минимальный контрпример, цикл, внутренние и внешние части, запрещённые
конфигурации. Геометрическая редукция леммы 14 (с. 44–45) остаётся
схематической; определения, предварительные шаги и извлечение двух
стандартных конфигураций расписаны явно. Теорема о стягиваемом ребре
и другое книжное доказательство здесь не используются.

\textbf{1. Минимальный контрпример и трёхсвязность.}
Выберем непланарный граф \(G\) без двух запрещённых подграфов,
минимальный сначала по \(n\), затем по \(m\).
Любой его собственный подграф планарен.
Граф связен и не имеет точек сочленения: иначе меньшие части планарны
и их укладки склеиваются по билету 16. При \(n\ge3\) мост также
давал бы точку сочленения. Малые графы планарны, поэтому они не контрпримеры.

Покажем, что нет и разделяющей пары \(\{a,b\}\).
Разобьём компоненты \(G-\{a,b\}\) на две непустые группы и добавим
к каждой \(a,b\), получив \(G_1,G_2\). Каждый компонент смежен и с
\(a\), и с \(b\), иначе одна из этих вершин была бы точкой сочленения.
Значит, в каждой части есть \(ab\)-путь с внутренними вершинами в этой части.
Рассмотрим \(G_i+ab\). Если в нём есть запрещённое подразбиение,
оно использует новое ребро \(ab\); заменим его путём через другую часть.
Получим запрещённое подразбиение в \(G\), противоречие.
По минимальности обе дополненные части планарны. Выносим \(ab\)
на внешнюю грань и склеиваем укладки по ребру с разных сторон;
если ребра не было, удаляем его. Получаем планарность \(G\), противоречие.
Таким образом, \(G\) трёхсвязен: удаление любых одной или двух вершин
сохраняет связность.

\textbf{2. Удалённое ребро и цикл (лемма 13 лекций).}
Выберем \(uv\in E(G)\), положим \(H=G-uv\). Граф \(H\) планарен
по минимальности и двусвязен. Действительно, для любой вершины \(z\)
граф \(G-z\) двусвязен по трёхсвязности \(G\), поэтому ребро \(uv\),
если оно осталось, не мост в \(G-z\). Значит, \(H-z\) связен.
По билету 5 в \(H\) есть простой цикл \(C\) через \(u,v\).

\begin{definition}[Части относительно цикла]
Частью относительно \(C\) называется либо хорда \(C\), либо компонента
\(H-V(C)\) вместе со всеми рёбрами к \(C\) и их концами на \(C\).
Эти концы — точки прикрепления. В плоской укладке внутренности каждой
части лежат по одну сторону \(C\). Обозначаем внешние части \(A_j\),
внутренние — \(B_i\), как в лекциях. Это \emph{не} мосты-рёбра из билета 4.
\end{definition}
На сфере можно выбирать укладку \(H\), цикл \(C\) через \(u,v\) и его
сторону так, чтобы на выбранной внутренней стороне было как можно больше
граней. Максимум существует: число граней конечно и фиксировано формулой Эйлера.
Именно такой экстремальный выбор позволяет перестраивать цикл,
сравнивая число захваченных граней (рисунок с. 42–43).

Если внешняя часть имеет хотя бы три точки прикрепления, внутри неё
можно выбрать дерево с тремя выходами на \(C\). Между тремя выходами
есть дуга \(C\), внутренность которой не содержит ни \(u\), ни \(v\).
Заменяя её путём по двум ветвям дерева, выбираем новую внутреннюю сторону,
содержащую прежнюю и ещё область между дугой и деревом. Одна из двух
остальных областей остаётся снаружи. Число внутренних граней растёт,
противоречие максимальности. В связном графе без точек сочленения у
части не меньше двух точек прикрепления. Поэтому у каждой \(A_j\)
ровно две. Внешняя часть не может иметь внутренних вершин: её две
точки прикрепления разделяли бы \(G\). Следовательно, все \(A_j\) — хорды.

\begin{center}
\begin{tikzpicture}[scale=.65,every node/.style={circle,fill=cbBlue,inner sep=2pt}]
\draw (0,0) circle (1.5);
\node (u) at (0,1.5) {};\node[fill=none,above] at (u) {$u$};
\node (v) at (0,-1.5) {};\node[fill=none,below] at (v) {$v$};
\node (a) at (1.5,0) {};\node[fill=none,right] at (a) {$a$};
\node (b) at (-1.5,0) {};\node[fill=none,left] at (b) {$b$};
\draw[cbOrange,thick] (a) .. controls (3,3.3) and (-3,3.3) .. (b);
\node[fill=none] at (1.6,2.1) {$A_j$};
\node (z) at (0,0) {};
\draw[cbGreen,thick] (z)--(u);\draw[cbGreen,thick] (z)--(v);
\draw[cbGreen,thick] (z)--(a);\draw[cbGreen,thick] (z)--(b);
\node[fill=none] at (.65,.5) {$B_i$};\node[fill=none] at (-1.1,1.2) {$C$};
\end{tikzpicture}
\end{center}
Внешняя хорда \(A_j\) и внутренняя часть \(B_i\) прикреплены к одному
циклу \(C\). Само ребро \(uv\) удалено; задача — вернуть его без пересечений.

\textbf{3. Разделение частей (лемма 14 лекций).}
Две пары точек \(a,b\) и \(c,d\) разделяют друг друга на цикле,
если четыре точки различны и встречаются в порядке \(a,c,b,d\).
Пути между такими парами не могут лежать по одну сторону \(C\):
путь первой пары вместе с дугой цикла отделяет концы второго по теореме Жордана.
Более общо, две части конфликтуют, если имеют чередующиеся пары точек
прикрепления или не менее трёх общих точек прикрепления.
Для хорды \(ab\) конфликт с частью \(B\) означает наличие у \(B\)
точек в обеих открытых дугах \(C-\{a,b\}\).

Ребро \(uv\) тоже нужно разместить по одну из сторон цикла.
В лекционной геометрической редукции сравниваются перестановки частей
между сторонами и замены дуг \(C\) путями через части. Если \(uv\)
разместить нельзя, выделяется конфликт внутренних и внешних частей,
разделяющий также \(u,v\); сокращение несущественных ветвей даёт
одну из запрещённых конфигураций ниже. Если конфликта нет, части
переносятся на подходящие стороны и дуга \(uv\) дорисовывается —
получается укладка \(G\), противоречие.

\textit{Статус этого шага.} Полный перебор взаимного расположения точек
прикрепления здесь, как и в лекционной схеме, не доказывается.
В частности, произвольный конфликт двух частей сам по себе ещё не
доказывает непланарность: требуется невозможность распределить все части
и \(uv\) по сторонам. Рисунки ниже объясняют заключительные конфигурации,
а не заменяют эту редукцию необоснованным утверждением о любом конфликте.

\textbf{4. Как возникают запрещённые подграфы.}
\emph{Конфигурация \(K_{3,3}\).} На цикле по порядку лежат
\(u,a,c,v,b,d\), а вне его выбраны три пути \(uv,ab,cd\)
с попарно непересекающимися внутренностями, не встречающими \(C\).
Сохраним шесть дуг между соседними отмеченными вершинами и эти три пути.
Это подразбиение \(K_{3,3}\) с долями \(\{u,c,b\}\), \(\{a,v,d\}\).
Каждая главная вершина соединена с тремя вершинами другой доли;
все девять путей имеют разные внутренние вершины.

\emph{Конфигурация \(K_5\).} На цикле по порядку лежат \(u,a,v,b\),
имеются противоположные пути \(uv,ab\), а внутренняя часть содержит
вершину \(z\) с четырьмя внутренне непересекающимися выходами к
\(u,a,v,b\). Сохраняем четыре дуги цикла, два противоположных пути
и четыре выхода: получаем десять путей подразбиения \(K_5\)
с главными вершинами \(u,a,v,b,z\).
В геометрическом разборе выбираются именно непересекающиеся ветви;
при нескольких развилках выделяются соответствующие пути и удаляется лишнее.

\begin{center}
\begin{tikzpicture}[scale=.8,every node/.style={circle,fill=cbBlue,inner sep=2pt}]
\foreach \i/\lab in {1/u,2/a,3/c,4/v,5/b,6/d}{
 \node (p\i) at ({90-60*(\i-1)}:1.25) {};
 \node[fill=none] at ({90-60*(\i-1)}:1.55) {$\lab$};}
\draw (p1)--(p2)--(p3)--(p4)--(p5)--(p6)--(p1);
\draw[cbOrange,thick] (p1)--(p4);\draw[cbGreen,thick] (p2)--(p5);
\draw[cbPurple,thick] (p3)--(p6);
\node[rectangle,fill=none] at (0,-2) {$K_{3,3}$};
\begin{scope}[xshift=5cm]
\foreach \i/\lab in {1/u,2/a,3/v,4/b}{
 \node (q\i) at ({90-90*(\i-1)}:1.25) {};
 \node[fill=none] at ({90-90*(\i-1)}:1.55) {$\lab$};}
\node (z) at (0,0) {};\node[fill=none,above right] at (z) {$z$};
\draw (q1)--(q2)--(q3)--(q4)--(q1);
\foreach \i in {1,...,4}{\draw (z)--(q\i);}
\draw[cbOrange,thick] (q1) to[bend left=28] (q3);
\draw[cbGreen,thick] (q2) to[bend left=35] (q4);
\node[rectangle,fill=none] at (0,-2) {$K_5$};
\end{scope}
\end{tikzpicture}
\end{center}
Это схемы подграфов, а не плоские укладки: пересечения линий без отмеченной
вершины не означают общей вершины соответствующих путей.

Итак, геометрический шаг либо даёт укладку \(G\), либо выделяет
запрещённое подразбиение. Оба исхода противоречат выбору контрпримера.
Это завершает \emph{схему} достаточности в форме лекций.

\endgroup

\ticket{Хроматический многочлен и рекуррентные формулы}
\label{ticket:18}
\begin{definition}[Правильная раскраска и хроматический многочлен]
Правильная раскраска в \(k\) цветов — функция \(\varphi:V\to\{1,\ldots,k\}\),
для которой \(\varphi(u)\ne\varphi(v)\) при \(uv\in E\).
Разрешены неиспользованные цвета; цвета различимы.
Число таких раскрасок обозначается \(\chi_G(k)\).
Хроматическое число \(\chi(G)\) — минимальное число цветов правильной раскраски.
\end{definition}
\begin{lemma}[Удаление и стягивание]
Для ребра \(e=uv\) простого графа
\[
 \boxed{\chi_G(k)=\chi_{G-e}(k)-\chi_{G/e}(k).}
\]
После стягивания кратные рёбра можно заменить одиночными.
\end{lemma}
\begin{proof}
Разобьём раскраски \(G-e\) на два класса.
При различных цветах \(u,v\) это ровно раскраски \(G\).
При одинаковых цветах \(u,v\) отождествляем их: получается биекция
с раскрасками \(G/e\). Кратные рёбра задают то же самое неравенство цветов,
поэтому их число на раскраски не влияет. Складывая мощности, получаем формулу.
\end{proof}
\begin{graphillustration}{Лекционный пример удаления и стягивания одного ребра квадрата. Слева сохраняются четыре вершины, справа концы $u,v$ отождествлены.}{fig:18-delete-contract}
\begin{tikzpicture}[font=\small]
\node[gi vertex] (u) at (0,1.6) {$u$};
\node[gi vertex] (v) at (0,0) {$v$};
\node[gi vertex] (a) at (1.6,0) {$a$};
\node[gi vertex] (b) at (1.6,1.6) {$b$};
\draw[gi edge] (v) -- (a);
\draw[gi edge] (a) -- (b);
\draw[gi edge] (b) -- (u);
\draw[gi removed] (u) -- (v);
\node[gi note] at (0.8,-0.8) {$G-e=P_4$};
\node[gi note] at (3,0.8) {$\longleftarrow$};
\begin{scope}[shift={(4.2,0)}]
\node[gi vertex] (u) at (0,1.6) {$u$};
\node[gi vertex] (v) at (0,0) {$v$};
\node[gi vertex] (a) at (1.6,0) {$a$};
\node[gi vertex] (b) at (1.6,1.6) {$b$};
\draw[gi emphasis] (u) -- (v);
\draw[gi edge] (v) -- (a);
\draw[gi edge] (a) -- (b);
\draw[gi edge] (b) -- (u);
\node[gi note] at (0.8,-0.8) {$G=C_4$, $e=uv$};
\end{scope}
\node[gi note] at (7.2,0.8) {$\longrightarrow$};
\begin{scope}[shift={(8.5,0)}]
\node[gi selected] (w) at (0,0.8) {$u=v$};
\node[gi vertex] (a) at (1.6,0) {$a$};
\node[gi vertex] (b) at (1.6,1.6) {$b$};
\draw[gi edge] (w) -- (a);
\draw[gi edge] (a) -- (b);
\draw[gi edge] (b) -- (w);
\node[gi note] at (0.8,-0.8) {$G/e=C_3$};
\end{scope}
\node[gi note] at (5,-1.8) {$\chi_{C_4}=\chi_{P_4}-\chi_{C_3}$};
\end{tikzpicture}
\end{graphillustration}


\begin{graphillustration}{Два класса раскрасок $G-e$: показано по одному примеру. Цифры — цвета; равенство цветов концов удалённого ребра соответствует стягиванию.}{fig:18-color-classes}
\begin{tikzpicture}[font=\small]
\begin{scope}[shift={(0,0)}]
\node[gi selected,label=left:$u$] (u) at (0,1.6) {1};
\node[gi selected,label=left:$v$] (v) at (0,0) {2};
\node[gi vertex,label=right:$a$] (a) at (1.6,0) {1};
\node[gi vertex,label=right:$b$] (b) at (1.6,1.6) {2};
\draw[gi edge] (v) -- (a);
\draw[gi edge] (a) -- (b);
\draw[gi edge] (b) -- (u);
\node[gi note] at (0.8,-0.9) {$\varphi(u)\ne\varphi(v)$: раскраски $G$};
\end{scope}
\begin{scope}[shift={(6,0)}]
\node[gi selected,label=left:$u$] (u) at (0,1.6) {1};
\node[gi selected,label=left:$v$] (v) at (0,0) {1};
\node[gi vertex,label=right:$a$] (a) at (1.6,0) {3};
\node[gi vertex,label=right:$b$] (b) at (1.6,1.6) {2};
\draw[gi edge] (v) -- (a);
\draw[gi edge] (a) -- (b);
\draw[gi edge] (b) -- (u);
\node[gi note] at (0.8,-0.9) {$\varphi(u)=\varphi(v)$: раскраски $G/e$};
\end{scope}
\end{tikzpicture}
\end{graphillustration}


\begin{cor}[Добавление и стягивание]
Для несмежных \(u,v\)
\[
 \chi_G(k)=\chi_{G+uv}(k)+\chi_{G/(u=v)}(k).
\]
\end{cor}
\begin{proof}
Снова разделяем раскраски по равенству цветов \(u,v\).
Разные цвета соответствуют \(G+uv\), одинаковые — отождествлению \(u,v\).
\end{proof}
\begin{theorem}[Полиномиальность]
Для графа без петель \(\chi_G(k)\) является унитарным многочленом
степени \(n\) с целыми коэффициентами. Для пустого множества вершин \(\chi_G=1\).
\end{theorem}
\begin{proof}
Индукция по \(n+m\); кратные рёбра предварительно удаляем.
Для графа без рёбер вершины красятся независимо, \(\chi_G(k)=k^n\).
При наличии ребра используем удаление–стягивание: первый многочлен
по индукции унитарен степени \(n\), второй — степени \(n-1\).
Их разность унитарна степени \(n\) и имеет целые коэффициенты.
Полученный многочлен совпадает с числом раскрасок для всех натуральных \(k\).
Он единственен: ненулевой многочлен не может иметь бесконечно много корней.
\end{proof}
\begin{lemma}[Произведение по компонентам]
Если компоненты графа — \(G_1,\ldots,G_c\), то
\(\chi_G(k)=\prod_{j=1}^c\chi_{G_j}(k)\).
\end{lemma}
\begin{proof}
Выбор раскраски каждой компоненты независим, поскольку рёбер между
компонентами нет. Правило произведения даёт равенство для натуральных \(k\),
а значит, и тождество многочленов.
\end{proof}
\begin{example}
\(\chi_{K_n}(k)=k(k-1)\cdots(k-n+1)\): последовательно выбираем различные
цвета для всех вершин. Для цикла \(C_3\) рекурсия даёт
\(\chi_{C_3}=k(k-1)^2-k(k-1)=k(k-1)(k-2)\).
Лекционный пример \(C_4\): удаление ребра даёт \(P_4\), стягивание — \(C_3\), поэтому
\[
\chi_{C_4}(t)=t(t-1)^3-t(t-1)(t-2)=t(t-1)(t^2-3t+3).
\]
При прямом подсчёте фиксируем цвет первой вершины. Если противоположная
имеет тот же цвет, получается \(t(t-1)^2\) раскрасок; если другой —
\(t(t-1)(t-2)^2\). Сумма даёт ту же формулу.
Петля запрещает правильные раскраски: \(\chi_G\equiv0\).
\end{example}

\ticket{Хроматический многочлен дерева}
\label{ticket:19}
\begin{definition}
Дерево — непустой связный граф без циклов. В этом билете граф простой.
\(\chi_G(k)\) — число правильных раскрасок различимыми цветами из \([k]\),
продолженное до многочлена по билету 18.
\end{definition}
\begin{lemma}[Удаление листа]
Если \(v\) — лист, то \(\chi_G(k)=(k-1)\chi_{G-v}(k)\).
\end{lemma}
\begin{proof}
Для каждой раскраски \(G-v\) цвет единственного соседа \(v\) запрещён,
а остальные \(k-1\) цветов разрешены. Продолжения независимы от других вершин.
Равенство чисел раскрасок для всех натуральных \(k\) даёт тождество многочленов.
\end{proof}
\begin{theorem}[Характеристика дерева]
Для простого графа на \(n\ge1\) вершинах
\[
 \boxed{G\text{ — дерево}\quad\Longleftrightarrow\quad
 \chi_G(k)=k(k-1)^{n-1}.}
\]
\end{theorem}
\begin{proof}
Необходимость: при \(n=1\) раскрасок \(k\). При \(n>1\) удаляем лист;
оставшийся граф — дерево. Индукция и лемма дают
\((k-1)\cdot k(k-1)^{n-2}=k(k-1)^{n-1}\).
Можно также выбрать корень: его цвет произволен, а каждая другая
вершина выбирает любой цвет, кроме цвета родителя. В дереве других
ограничений между уже выбранными предками нет.

Достаточность: степень многочлена равна числу вершин по билету 18.
Коэффициент при \(k\) у \(k(k-1)^{n-1}\) равен \((-1)^{n-1}\ne0\).
По лекционной теореме билета 20 наименьшая ненулевая степень равна
числу компонент (её доказательство не использует данный билет).
Значит, \(G\) связен.
Для простого графа коэффициент при \(k^{n-1}\) равен \(-m\)
(доказательство дано в билете 20 и не использует данный билет).
Сравнение с разложением \(k(k-1)^{n-1}\) даёт \(m=n-1\)
при \(n\ge2\). Связный граф с \(n-1\) рёбрами — дерево по билету 6.
При \(n=1\) простой граф сам является одноточечным деревом.
\end{proof}
\begin{cor}[Многочлен леса]
Для леса с \(n\) вершинами и \(c\) компонентами
\[
 \chi_G(k)=k^c(k-1)^{n-c}.
\]
\end{cor}
\begin{proof}
Каждая компонента — дерево. Перемножаем её формулу;
сумма чисел вершин равна \(n\), число множителей \(k\) равно \(c\).
Для пустого леса получаем \(1\).
\end{proof}
\textbf{Зачем нужна простота.} Если разрешить кратные рёбра, удвоение
любого ребра дерева не меняет хроматический многочлен, но полученный
мультиграф уже не дерево в обычном смысле. Для мультиграфов формула
характеризует дерево, полученное после удаления кратностей, при отсутствии петель.

\ticket{Коэффициенты хроматического многочлена}
\label{ticket:20}
\source{Лекции, с. 50–51; Карпов, §14.1; пропущенные шаги индукции раскрыты составителем}
\begin{definition}
Пусть \(G\) простой, \(n\) — число вершин, \(m\) — число рёбер,
\(c=c(G)\). Запишем
\[
 \chi_G(k)=\sum_{j=0}^n(-1)^j b_j(G)k^{n-j}.
\]
Знакопеременность означает \(b_j\ge0\); ниже уточняется, какие из них строго положительны.
\end{definition}
\begin{theorem}[Старшие коэффициенты]
Старший коэффициент \(\chi_G\) равен 1, а при \(n\ge1\) коэффициент
при \(k^{n-1}\) равен \(-m\).
\end{theorem}
\begin{proof}
Индукция по числу рёбер, одновременно для всех чисел вершин.
При \(m=0\) имеем \(\chi_G(k)=k^n\); второй коэффициент нулевой.
Пусть \(e\in E\). По индукции
\[
\chi_{G-e}(k)=k^n-(m-1)k^{n-1}+\text{члены меньшей степени},
\]
а \(\chi_{G/e}\) унитарен степени \(n-1\) по билету 18.
Из \(\chi_G=\chi_{G-e}-\chi_{G/e}\) получаем старший коэффициент 1
и следующий \(-(m-1)-1=-m\). Кратные рёбра после стягивания
убираются, что не влияет на унитарность; сравнивать их число здесь не требуется.
При \(n=1\) имеем \(m=0\) и \(\chi_G=k\).
\end{proof}
\begin{lemma}[Рекурсия для модулей коэффициентов]
Для ребра \(e\)
\[
 b_j(G)=b_j(G-e)+b_{j-1}(G/e),
\]
где коэффициенты с индексами вне допустимого диапазона считаются нулевыми,
а кратные рёбра после стягивания удаляются.
\end{lemma}
\begin{proof}
Подставляем разложения многочленов в
\(\chi_G=\chi_{G-e}-\chi_{G/e}\). У второго графа на одну вершину меньше,
поэтому его коэффициент при \(k^{n-j}\) равен
\((-1)^{j-1}b_{j-1}(G/e)\). Вычитание меняет знак на \((-1)^j\),
и после сравнения коэффициентов получается сумма.
\end{proof}
\begin{theorem}[Чередование знаков]
\[
 b_j(G)>0\quad(0\le j\le n-c),\qquad
 b_j(G)=0\quad(j>n-c).
\]
Таким образом, ненулевые коэффициенты при степенях
\(k^n,k^{n-1},\ldots,k^c\) строго чередуются по знаку,
а коэффициенты более низких степеней нулевые.
\end{theorem}
\begin{proof}
Индукция по \(n+m\). Для графа без рёбер \(\chi_G=k^n\), \(c=n\),
так что положителен только \(b_0=1\).
Пусть есть ребро \(e\), положим \(r=n-c\ge1\).
Стягивание не меняет число компонент, следовательно,
\((n-1)-c(G/e)=r-1\). При удалении немоста число компонент не меняется,
а при удалении моста увеличивается на 1.
По индукции оба слагаемых рекурсии неотрицательны.
Если \(e\) не мост, первое слагаемое положительно для \(0\le j\le r\)
и равно нулю выше \(r\); второе может быть положительно только при
\(1\le j\le r\). Если \(e\) мост, первое положительно при
\(0\le j\le r-1\), второе при \(1\le j\le r\).
В обоих случаях сумма положительна ровно для \(0\le j\le r\).
\end{proof}
\begin{example}
Для дерева на четырёх вершинах
\(\chi_G(k)=k(k-1)^3=k^4-3k^3+3k^2-k\).
Старший коэффициент 1, второй \(-3=-m\), ненулевые знаки чередуются;
свободный член равен нулю.
\end{example}

\begin{cor}[Наименьшая ненулевая степень]
Наименьшая степень с ненулевым коэффициентом в \(\chi_G(k)\) равна \(c(G)\).
В частности, непустой граф связен тогда и только тогда,
когда коэффициент при \(k\) ненулевой.
\end{cor}
\begin{proof}
В теореме о чередовании последний положительный модуль — \(b_{n-c}\),
стоящий при \(k^c\); все следующие коэффициенты нулевые.
\end{proof}
\textbf{Дополнительное доказательство по включениям–исключениям.}
Следующая формула отсутствует в скане и не нужна для лекционной индукции.
\begin{lemma}[Включения–исключения]
\[
 \chi_G(k)=\sum_{S\subseteq E}(-1)^{|S|}k^{c(V,S)},
\]
где \(c(V,S)\) — число компонент остовного подграфа с рёбрами \(S\).
\end{lemma}
\begin{proof}
Из всех \(k^n\) раскрасок исключаем события \(B_e\): концы ребра \(e\)
имеют одинаковый цвет. При выполнении всех событий для \(e\in S\)
вершины каждой компоненты \((V,S)\) должны быть одноцветны; цвета
разных компонент выбираются независимо, поэтому пересечение событий
имеет мощность \(k^{c(V,S)}\). Формула включений–исключений даёт нужную сумму.
Для полноты: раскраска с \(t\) нарушенными рёбрами получает в этой сумме
вес \(\sum_{r=0}^t(-1)^r\binom tr=(1-1)^t\), равный 1 при \(t=0\)
и 0 при \(t>0\), что непосредственно доказывает формулу.
\end{proof}

\ticket{Теорема Брукса}
\label{ticket:21}
\source{Лекции, с. 55–60; Карпов, §4.1–4.2; детали лекционного доказательства Ловаса восстановлены}
\begin{definition}[Раскраска]
\(\chi(G)\) — минимальное число цветов правильной раскраски вершин;
\(\Delta=\Delta(G)\) — максимальная степень. В этом билете граф простой и связный.
Жадный алгоритм всегда даёт \(\chi(G)\le\Delta+1\): в любом порядке
каждая очередная вершина видит не более \(\Delta\) уже запрещённых цветов.
\end{definition}
Граф называется \textbf{регулярным}, если все его вершины имеют одну степень.
Если эта степень равна \(d\), граф называют \(d\)-регулярным.\par\medskip
\begin{lemma}[Вершина неполной степени]\label{lem:greedy}
Если \(G\) связен, \(\Delta(G)\le d\), \(d\ge1\), и есть вершина \(r\)
степени меньше \(d\), то \(\chi(G)\le d\).
\end{lemma}
\begin{proof}
Выберем остовное дерево с корнем \(r\) и красим вершины в порядке
убывания расстояния до корня (порядок на одном уровне произволен).
У каждой \(x\ne r\) родитель ещё не окрашен, поэтому окрашенных
соседей не больше \(d-1\). Из \(d\) цветов всегда есть свободный.
Корень красим последним; у него всего меньше \(d\) соседей,
так что свободный цвет тоже найдётся.
\end{proof}
\begin{lemma}[Тройка Ловаса]\label{lem:lovasz}
Если граф двусвязен, неполон и \(\delta(G)\ge3\), существуют вершины
\(a,b,c\), для которых \(ab,bc\in E\), \(ac\notin E\),
а \(G-\{a,c\}\) связен.
\end{lemma}
\begin{proof}
Если граф трёхсвязен (имеет не менее четырёх вершин и остаётся связным
после удаления любых одной или двух), возьмём кратчайший путь между
двумя несмежными вершинами. Его первые три вершины \(a,b,c\)
удовлетворяют \(ac\notin E\), иначе путь не кратчайший.
Связность после удаления \(a,c\) следует из трёхсвязности.

Если граф не трёхсвязен, существует разделяющая пара \(\{b,z\}\).
Граф \(G'=G-b\) связен по двусвязности \(G\), но имеет точку сочленения \(z\).
По дереву блоков (билет 5) у \(G'\) есть два разных крайних блока
\(B_a,B_c\). Обозначим единственные точки их присоединения к остальному
графу через \(s_a,s_c\) (они могут совпасть).
Вершина \(b\) смежна с некоторой внутренней вершиной каждого крайнего
блока. Иначе удаление \(s_a\), соответственно \(s_c\), отделяло бы
внутренние вершины блока и в исходном графе \(G\), вопреки двусвязности.
Выберем таких соседей \(a\in B_a\setminus\{s_a\}\),
\(c\in B_c\setminus\{s_c\}\).
Ребра \(ac\) нет, поскольку внутренние вершины разных крайних блоков
не соединены рёбрами.

Крайний блок не может состоять из одного ребра: его внутренняя вершина
имела бы степень 1 в \(G'\) и не больше 2 в \(G\), вопреки \(\delta(G)\ge3\).
Следовательно, оба выбранных блока невырожденно двусвязны.
Удаление \(a\) из \(B_a\) и \(c\) из \(B_c\) не нарушает их связности;
точки присоединения остаются, и потому \(G'-\{a,c\}\) связен.
У \(b\) не менее трёх соседей, так что после удаления двух вершин
хотя бы один сосед остаётся. Добавление \(b\) сохраняет связность.
\end{proof}
\begin{theorem}[Брукс]
Если связный простой граф не является ни полным графом,
ни циклом нечётной длины, то
\[
 \boxed{\chi(G)\le\Delta(G).}
\]
Исключения удовлетворяют \(\chi(K_n)=n=\Delta+1\),
\(\chi(C_{2r+1})=3=\Delta+1\).
\end{theorem}
\begin{proof}
При \(\Delta\le1\) связный граф — \(K_1\) или \(K_2\), оба исключены.
При \(\Delta=2\) связный граф — путь либо цикл:
если есть вершина степени 1, двигаясь из неё по единственному продолжению,
проходим весь граф, иначе возникла бы несвязность; если все степени 2,
тот же обход замыкается в цикл и исчерпывает граф.
Путь и чётный цикл красятся в два цвета поочерёдно;
нечётный цикл исключён.

Пусть \(d=\Delta\ge3\). Если хотя бы одна вершина имеет степень меньше
\(d\), применяем лемму \ref{lem:greedy}. Остаётся \(d\)-регулярный граф.
Если есть точка сочленения \(v\), обозначим компоненты \(G-v\) через
\(H_1,\ldots,H_t\), \(t\ge2\). Каждый граф \(G_i=G[V(H_i)\cup\{v\}]\)
связен, имеет максимальную степень не больше \(d\), а степень \(v\)
в нём меньше \(d\), поскольку хотя бы один сосед \(v\) находится
в другой компоненте. По первой лемме каждый \(G_i\) красится
в \(d\) цветов. Перестановкой названий цветов делаем цвет \(v\)
одинаковым во всех раскрасках и объединяем их. Между разными \(H_i\)
рёбер нет, поэтому раскраска правильна.

Наконец, пусть \(G\) двусвязен. По лемме \ref{lem:lovasz} есть
\(a,b,c\) с несмежными \(a,c\), общим соседом \(b\) и связным
\(G-\{a,c\}\). Красим \(a,c\) в один цвет. В оставшемся графе
берём остовное дерево с корнем \(b\) и красим снизу вверх.
У каждой вершины, кроме \(b\), остаётся неокрашенный родитель,
поэтому запрещено не более \(d-1\) цветов, и выбор возможен.
Корень \(b\) видит \(d\) соседей, но два из них, \(a,c\), имеют
одинаковый цвет. Следовательно, занято не более \(d-1\) цветов;
свободным цветом красим \(b\). Получена правильная \(d\)-раскраска.

В полном графе все вершины попарно смежны, поэтому нужны \(n\) цветов.
Нечётный цикл нельзя раскрасить в два цвета из-за чередования цветов
вдоль обхода; трёх цветов достаточно, если одну вершину покрасить
третьим цветом, а оставшийся путь — двумя. Исключения проверены.
\end{proof}
Для несвязного графа утверждение применяют к каждой компоненте отдельно;
хроматическое число всего графа — максимум чисел его компонент, поскольку
между ними нет ограничений на повторное использование цветов.

\end{document}
